\documentclass[11pt,a4paper]{article}
\usepackage[margin=27mm,headheight=15pt]{geometry}
\usepackage{fontspec}
\setmainfont{Latin Modern Roman}
\usepackage{amsmath,amssymb,amsthm,mathtools,mathrsfs,enumitem}
\usepackage{xcolor,hyperref,bookmark,fancyhdr,array,longtable,needspace,etoolbox}
\definecolor{linkblue}{HTML}{164B76}
\hypersetup{colorlinks=true,linkcolor=linkblue,citecolor=linkblue,urlcolor=linkblue,
 pdftitle={Plurigenera and canonical algebras in analytic deformations},pdfauthor={Tseng}}
\pagestyle{fancy}\fancyhf{}
\fancyhead[L]{\small Plurigenera and canonical algebras}
\fancyhead[R]{\small Tseng}\fancyfoot[C]{\thepage}
\setlength{\parskip}{3pt}\setlength{\parindent}{1.2em}
\setlist{itemsep=4pt,topsep=5pt}
\emergencystretch=2em
\clubpenalty=10000
\widowpenalty=10000
\displaywidowpenalty=10000
\numberwithin{equation}{section}
\newtheorem{theorem}{Theorem}[section]
\newtheorem{proposition}[theorem]{Proposition}
\newtheorem{lemma}[theorem]{Lemma}
\newtheorem{corollary}[theorem]{Corollary}
\newtheorem{maintheorem}{Theorem}
\renewcommand{\themaintheorem}{\Alph{maintheorem}}
\theoremstyle{definition}
\newtheorem{definition}[theorem]{Definition}
\theoremstyle{remark}
\newtheorem{remark}[theorem]{Remark}
\newtheorem{example}[theorem]{Example}
\newcommand{\Q}{\mathbb Q}\newcommand{\R}{\mathbb R}
\newcommand{\C}{\mathbb C}\newcommand{\Z}{\mathbb Z}
\newcommand{\PP}{\mathbf P}\newcommand{\HH}{\mathbb H}
\newcommand{\OO}{\mathcal O}\newcommand{\cX}{\mathcal X}
\newcommand{\cY}{\mathcal Y}\newcommand{\cB}{\mathcal B}
\newcommand{\cA}{\mathcal A}\newcommand{\cF}{\mathcal F}
\newcommand{\cL}{\mathcal L}\newcommand{\cC}{\mathcal C}
\newcommand{\Supp}{\operatorname{Supp}}\newcommand{\Sing}{\operatorname{Sing}}
\newcommand{\ord}{\operatorname{ord}}\newcommand{\per}{\operatorname{per}}
\newcommand{\Spec}{\operatorname{Spec}}\newcommand{\vol}{\operatorname{vol}}
\newcommand{\Proj}{\operatorname{Proj}}
\newcommand{\rank}{\operatorname{rank}}\newcommand{\dd}{\mathrm d}
\newcommand{\id}{\operatorname{id}}\newcommand{\lct}{\operatorname{lct}}
\newcommand{\BC}{\mathrm{BC}}\newcommand{\reg}{\mathrm{reg}}
\title{Plurigenera and canonical algebras in analytic deformations}
\author{Tseng}
\date{11 November 2026}
\begin{document}
\maketitle
\begin{abstract}
We prove invariance of all reflexive plurigenera in proper flat
analytic families whose central fiber is compact Kähler and has
canonical singularities. The base may be singular or nonreduced,
and a single neighborhood suffices for every degree. This proves Kollár's Conjecture~16 with a compact
Kähler central fiber and implies Siu's conjecture for smooth families.
We also prove invariance of rounded logarithmic plurigenera for smooth
families of compact Kähler manifolds with a relative simple normal
crossings boundary and real coefficients in $[0,1)$.

We adapt Levine's infinitesimal extension method to klt central pairs
whose log canonical divisor is semiample and whose underlying space
admits a compact Kähler modification. An argument with finitely many
coherent sheaves gives a common neighborhood and relative semiampleness.
The birational reduction uses OpenAI's logarithmic subadditivity,
abundance and model theorems. Bingener's deformation theorem and a
disk criterion give the passage to arbitrary analytic bases.
Over a disk, we prove finite generation, flatness and base change for
the relative canonical algebra when the central fiber is canonical
and is either compact Kähler or admits a compact Kähler modification
and has semiample canonical divisor. Its relative Proj is normal,
flat and projective. For smooth disk families whose central fiber
belongs to class~$\mathcal C$ and has semiample canonical bundle,
all pluricanonical cohomology dimensions are constant on one
neighborhood, simultaneously in every integer tensor degree.
\end{abstract}
\begingroup
\small\noindent
\textit{2020 Mathematics Subject Classification.}
Primary 32G05, 32J27; Secondary 14E30.\par
\noindent\textit{Keywords.}
Plurigenera, deformation invariance, compact Kähler spaces,
canonical singularities, relative canonical algebras.
\endgroup
\setcounter{tocdepth}{1}
\begingroup
\small
\setlength{\parskip}{0pt}
\makeatletter
\patchcmd{\l@section}{1.0em}{0.30em}{}{\errmessage{Unable to set contents spacing}}
\makeatother
\tableofcontents
\endgroup
\clearpage
\section{Introduction and main results}\label{sec:introduction}

For a normal compact complex space with canonical singularities, write
\[
 \omega_X^{[m]}=(\omega_X^{\otimes m})^{**},\qquad
 P_m(X)=h^0(X,\omega_X^{[m]})\quad(m\ge1).
\]
The deformation invariance of these dimensions is a basic question in
complex geometry. Iitaka proved invariance for compact complex surfaces
\cite{Iitaka1970}. Siu established it for smooth projective families,
first in general type and then in arbitrary Kodaira dimension
\cite{Siu1998,Siu2002}, and formulated the corresponding conjecture for
families of compact Kähler manifolds
\cite[Conjecture~0.4]{Siu2002}. Kawamata gave an algebraic proof in
general type and proved deformation stability of canonical singularities
\cite{Kawamata1999}; Takayama treated projective families with canonical
fibers \cite{Takayama2007}. Păun subsequently simplified Siu's extension
argument \cite{Paun2007}.

For analytic families, Rao--Tsai prove invariance when all fibers are
smooth Moishezon manifolds \cite[Theorem~1.2\textup{(i)}]{RaoTsai2022}
and formulate a conjecture for canonical spaces in Fujiki's class
$\mathcal C$ \cite[Conjecture~1.10]{RaoTsai2022}. Kollár proves
invariance near a projective canonical central fiber of general type
\cite[Theorem~1]{Kollar2021}. We establish the central-fiber theorem
below for compact Kähler spaces in every Kodaira dimension. The proof
combines a central extension theorem with the logarithmic subadditivity,
abundance and birational model results of OpenAI \cite{OAI033,OAI034}.

\begin{maintheorem}[Canonical central fibers]\label{thm:canonical-central}
Let $f:\mathcal X\to S$ be a proper flat morphism of complex
analytic spaces, with $S$ not necessarily reduced, and let $0\in S$.
Suppose that $X_0$ is normal, connected and compact Kähler, with
canonical singularities. There is an open neighborhood $U$ of $0$
such that all fibers over $U$ have canonical singularities and
\[
 P_m(X_s)=P_m(X_0)\qquad(s\in U,\ m\ge1).
\]
If $S=\Delta$, the neighborhood $U$ may be chosen to be a disk
such that $\mathcal X_U=f^{-1}(U)$ has canonical singularities.
In this case, for the restricted family $f_U:\mathcal X_U\to U$
and every $m\ge1$, the sheaf
$(f_U)_*\OO_{\mathcal X_U}(mK_{\mathcal X_U/U})$ is locally free
and commutes with holomorphic base change, and every section of
$\omega_{X_0}^{[m]}$ extends over $\mathcal X_U$.
\end{maintheorem}

For a projective central fiber, this proves Kollár's
Conjecture~16 \cite{Kollar2021}, including Conjecture~17 as a
special case. The Kähler hypothesis concerns only the central fiber.
The following consequence gives Siu's conjecture for smooth families.

\Needspace{10\baselineskip}
\begin{corollary}[Siu's conjecture]\label{cor:siu}\label{cor:canonical}
Let $f:\mathcal X\to S$ be a proper holomorphic submersion of complex manifolds with
connected compact Kähler fibers, where $S$ is a connected complex
manifold. Then $P_m(X_s)$ is independent of $s\in S$ for every
$m\ge1$. The sheaves $f_*K_{\mathcal X/S}^{\otimes m}$ are locally
free and commute with holomorphic base change.
\end{corollary}

\begin{proof}
Theorem~\ref{thm:canonical-central} makes each $P_m(X_s)$ locally
constant on $S$, hence constant by connectedness.
Grauert's theorem \cite{Grauert1960} gives local freeness and
holomorphic base change.
\end{proof}

For $S=\Delta$, this is \cite[Conjecture~0.4]{Siu2002}.
The twistor family in Section~\ref{sec:example} illustrates the
fiberwise Kähler hypothesis: its total space is non-Kähler over
every nonempty open subset of the base.

\subsection{Logarithmic plurigenera}

For pairs, Hacon--McKernan--Xu prove invariance of rounded logarithmic
plurigenera in projective log-smooth families with coefficients in
$[0,1)$ \cite[Theorem~4.2]{HMX2018}. Partial results for ordinary
plurigenera in Kähler families include the smooth threefold case
\cite[Theorem~1.8]{HLR2026}, extension under nefness hypotheses
\cite{DNWZ2026,CRW2026}, canonical threefold families with a
$\Q$-factorial Kähler total space, and smooth Kähler families of
maximal Albanese dimension \cite[\mbox{Theorems}~1.4--1.5]{FLR2026}.
Li--Păun--Wang--Zhou prove the boundary-free result for families with
Kähler total space \cite[Theorem~1.2]{LPWZ2026}. In the same ambient
setting, Liu--Wan prove logarithmic extension and invariance, including
real coefficients by degreewise rounding
\cite[Theorem~1.1 and Corollary~1.2]{LW2026}.
Here the fibers may vary in a non-Kähler total space.

\Needspace{17\baselineskip}
\begin{maintheorem}[Logarithmic plurigenera]\label{thm:main}
Let $f:\mathcal X\to S$ be a proper holomorphic submersion of complex
manifolds, with $S$ connected and every fiber compact Kähler. Let
\[
 B=\sum_{i=1}^r b_iD_i,\qquad b_i\in\mathbb R,\quad 0\le b_i<1,
\]
where $\sum_iD_i$ is a relative simple normal crossings divisor.
For $m\ge1$, set
\[
 \mathcal L_m=\OO_{\mathcal X}
       \bigl(mK_{\mathcal X/S}+\lfloor mB\rfloor\bigr).
\]
Then $h^0(X_s,\mathcal L_m|_{X_s})$ is independent of $s\in S$.
For each $s\in S$, every section of $\mathcal L_m|_{X_s}$ extends
to a section of $\mathcal L_m$ on $f^{-1}(U)$ for some open
neighborhood $U$ of $s$. The sheaf $f_*\mathcal L_m$ is locally free
and commutes with holomorphic base change. If $S$ is Stein, every
fiber section extends over $\mathcal X$.
\end{maintheorem}

The case $B=0$ gives another proof of Corollary~\ref{cor:siu}.
For a fixed $m$, the rational boundary $m^{-1}\lfloor mB\rfloor$
reduces the logarithmic statement to an integral adjoint in that
degree, following \cite[proof of Theorem~4.2]{HMX2018}.

\subsection{Extension from a semiample central pair}

A \emph{compact Kähler modification} of a normal compact space $X$
is a proper bimeromorphic morphism $\widetilde X\to X$ from a
compact Kähler manifold. For smooth $X$, its existence characterizes
Fujiki's class $\mathcal C$. The extension argument applies in this
setting and allows singular klt boundaries.

\Needspace{16\baselineskip}
\begin{maintheorem}[Extension at a semiample central pair]\label{thm:endpoint}
Let $g:\mathcal Y\to\Delta$ be proper and flat, with normal connected
central fiber $X$, and suppose that $\mathcal Y$ is klt near $X$.
Let $\mathcal B$ be an effective rational horizontal divisor with
coefficients in $[0,1)$ such that
\[
 \mathcal A=K_{\mathcal Y/\Delta}+\mathcal B
\]
is $\mathbb Q$-Cartier. Suppose that adjunction identifies
$\mathcal A|_X$ with $A=K_X+B$, where $(X,B)$ is klt, $X$ admits a
compact Kähler modification, and $A$ is semiample. For each $m\ge1$
such that $m\mathcal B$ and $mB$ are integral, there is a disk $U_m$
about $0$ such that every section of $\OO_X(mA)$ extends to a
section of $\OO_{\mathcal Y}(m\mathcal A)$ over $g^{-1}(U_m)$.
\end{maintheorem}

The degree-one construction gives the following companion statement.

\begin{maintheorem}[Extension of canonical sections]\label{thm:degreeone}
Let $g:\mathcal Y\to\Delta$ be proper and flat, with normal connected
central fiber $X$. Suppose that $\mathcal Y$ is klt near $X$, that
$X$ is klt, and that $X$ admits a compact Kähler modification.
There is a disk $U$ about $0$ such that every section of $\omega_X$
extends to a section of $\OO_{\mathcal Y}(K_{\mathcal Y/\Delta})$
over $g^{-1}(U)$.
\end{maintheorem}

Levine developed the infinitesimal extension method for smooth
class-$\mathcal C$ central fibers carrying a smooth reduced
pluricanonical divisor \cite[Proposition~1.6 and Corollary~1.8]{Levine1983};
see also \cite{Levine1985}. The lower-degree argument in his
Corollary~1.8 is credited to N.~I.~Shepherd-Barron. Levine's
semiample corollary concerns families whose fibers are all in
class $\mathcal C$ \cite[Corollary~1.10]{Levine1983}, while his
torsion-canonical corollary assumes this only for the central fiber
\cite[Corollary~1.11]{Levine1983}. Kollár--Mori extend the method
to terminal singularities \cite[Theorem~12.5.5]{KM1992};
Kollár discusses its class-$\mathcal C$ scope and Nakayama's
projective extension results for canonical singularities in
\cite[Comments following Conjecture~17]{Kollar2021}; see
\cite[VI.3.15--16]{Nakayama2004}.

Theorem~\ref{thm:endpoint} extends this method to klt central pairs
and their admissible reflexive adjoint powers. Its canonical
specialization is Corollary~\ref{cor:canonical-endpoint}.
The uniformity argument of Lemma~\ref{lem:uniform} then gives
Corollary~\ref{cor:uniform-endpoint}: for a canonical central fiber
with a compact Kähler modification and semiample canonical divisor,
all plurigenera are constant on one disk and the relative canonical
divisor is semiample.

\subsection{Canonical algebras and higher cohomology}

The same uniformity argument also controls multiplication. It yields
a structural form of invariance in both canonical settings.

\begin{maintheorem}[Relative canonical algebras]
\label{thm:canonical-algebra}\label{cor:relative-canonical-algebra}
Let $f:\mathcal X\to\Delta$ be proper and flat, with normal connected
canonical central fiber $X_0$. Assume that either
\begin{enumerate}[label=\textup{(\alph*)}]
\item $X_0$ is compact Kähler; or
\item $X_0$ admits a compact Kähler modification and $K_{X_0}$ is semiample.
\end{enumerate}
There is a disk $U$ about $0$ on which the total space and all fibers
have canonical singularities. For the restricted family, set
\[
 \mathcal F_m=\OO_{\mathcal X}(mK_{\mathcal X/U}),\qquad
 \mathcal R(f)=\bigoplus_{m\ge0}f_*\mathcal F_m.
\]
The disk $U$ can be chosen so that the following hold.
\begin{enumerate}[label=\textup{(\roman*)}]
\item The $\OO_U$-algebra $\mathcal R(f)$ is finitely generated and
flat, and its graded pieces are locally free. For each holomorphic
map $u:T\to U$, there is a natural isomorphism of graded algebras
\[
 u^*\mathcal R(f)\simeq
 \bigoplus_{m\ge0}(f_T)_*(p^*\mathcal F_m),
\]
where $f_T:\mathcal X\times_U T\to T$ and
$p:\mathcal X\times_U T\to\mathcal X$ are the projections.
\item The relative analytic space $\mathcal C=\Proj_U\mathcal R(f)$
is normal, projective and flat over $U$. Its formation commutes
with holomorphic base change, and
\[
 \mathcal C_t\simeq\Proj\bigoplus_{m\ge0}
 H^0(X_t,\omega_{X_t}^{[m]}).
\]
\item If $\kappa(X_0)\ge0$, there is a proper flat family
$g:\mathcal Y\to U$ with canonical total space and canonical fibers,
and a bimeromorphic correspondence over $U$ inducing a bimeromorphic
map $X_t\dashrightarrow Y_t$ for every $t\in U$, such that
$K_{\mathcal Y/U}$ is relatively semiample. Under \textup{(b)}, one
can take $\mathcal Y=\mathcal X$.
\end{enumerate}
\end{maintheorem}

For a semiample logarithmic central pair, the corresponding adjoint
algebra has the same finite-generation and base-change properties
(Corollaries~\ref{cor:uniform-log-endpoint} and
\ref{cor:relative-adjoint-algebra}). For a smooth family whose central
fiber belongs to class~$\mathcal C$ and has semiample canonical bundle,
Corollary~\ref{cor:smooth-semiample-higher} proves constancy of all
$h^i(X_t,K_{X_t}^{\otimes m})$, for $i\ge0$ and $m\in\Z$, on one
common disk. The fixed-degree argument uses cyclic covers and Hodge
theory; Lemma~\ref{lem:uniform-higher} supplies the common disk.

\subsection{Proof strategy}

The argument over a disk has two parts: extension at a semiample central pair and
reduction to that pair by a finite birational program.

For extension, the cyclic-cover construction of Esnault--Viehweg
\cite[Section~3]{EV1992} identifies a character summand of
logarithmic top forms with the required reflexive adjoint sheaf.
Deligne's logarithmic degeneration theorem, in the compact
Kähler form of \cite[Theorem~3.2]{LRW2019}, gives an injective
central period map.
Its finite-order versions detect the lifting obstructions, which
vanish by Levine's local differential calculation. Proper coherence
then produces holomorphic extensions. All the Hodge-theoretic
constructions take place on a Kähler modification of the central
fiber. For smooth central fibers, a related infinitesimal approach
is developed by Cao--Păun, whose extension criteria use central
Kählerness and a weighted $L^2$ condition on the obstruction
\cite[Theorems~1.1--1.2]{CaoPaun2023}.

The reduction through minimal models and abundance originates in
Nakayama's work \cite{Nakayama1986}. For the disk case of
Theorem~\ref{thm:canonical-central}, we construct a finite negative
program beginning on the prescribed canonical central space. Its
steps include the mixed contractions of
\cite[Definition~7]{Kollar2021}. We use the analytic cone and
contraction theorems of Hacon--Xie \cite{HX2026}, relative finite
generation \cite{Fujino2022}, and OpenAI's abundance and birational
model theorems \cite{OAI034}. The termination proof develops the
affine-wall argument in \cite[proof of Theorem~3.13]{OAI034}.
Contraction extension and Nakayama's theorem deform this finite
program. At its semiample endpoint, grouping the canonical powers
modulo a generated Cartier multiple gives finitely many coherent
sheaves on relative projective space. Their uniform regularity
yields one neighborhood for all degrees and controls the canonical
algebra.

The passage to an arbitrary analytic base uses deformation openness.
Canonical singularities remain canonical \cite[VI.5.2]{Nakayama2004},
and Bingener's theorem ensures that the nearby fibers remain Kähler
\cite{Bingener1983,Kirschner2012}. The disk result therefore applies
at every point of one fixed neighborhood. A disk criterion, as in
\cite[paragraph~13]{Kollar2021}, then gives constancy of the entire
sequence of plurigenera: coordinate disks establish it on a resolution
of the base, and properness descends it to the original base.

For Theorem~\ref{thm:main}, we use the fixed-degree preparation in
\cite[proof of Theorem~4.2]{HMX2018}. Relative crepant stratum
blowups make the total and fiber pairs terminal while preserving
the chosen boundary denominator; the canonical preparation in
\cite[Lemma~4.2]{LW2026} is a related construction. We then subtract
the fixed boundary in the prescribed degree. This preserves its
integrality and excludes boundary components from the contracted
loci; it is the fixed-degree counterpart of the divisorial
negative-part construction in \cite[Lemma~4.3]{LW2026}.
The central log program deforms by \cite[Theorem~7.7]{HX2026}, and
Theorem~\ref{thm:endpoint} extends the resulting adjoint section.

Sections~\ref{sec:sheaves}--\ref{sec:uniform} prove the extension
and uniformity results. Sections~\ref{sec:canonical-steps}--\ref{sec:kollar-family}
prove Theorem~\ref{thm:canonical-central}, and
Section~\ref{sec:consequences} proves
Theorem~\ref{thm:canonical-algebra} and its cohomological and numerical
consequences. Sections~\ref{sec:finite}--\ref{sec:main-proof} establish
Theorem~\ref{thm:main}. Section~\ref{sec:example} gives the twistor
examples, and Appendix~\ref{app:dependencies} records the abundance
and nef-comparison results used in the birational arguments.

\section{Reflexive adjoints and birational section comparison}
\label{sec:sheaves}

\subsection{Conventions}

We work with complex analytic spaces and holomorphic maps. In
statements over a general analytic base, the base and total space
may be nonreduced. Holomorphic base changes may also be nonreduced.
Elsewhere spaces are reduced, except for the Artinian thickenings
used in the extension argument.
A normal connected space is irreducible. Families over
a disk are considered near the central fiber; shrinking the base includes
restricting the total space. We write $t$ for the disk coordinate and use
$\mathrm dt$ to identify relative and absolute canonical sheaves in
meromorphic section comparisons.

For a smooth family, a divisor has \emph{relative simple normal
crossings} (relative SNC) if it is locally defined by
$z_1\cdots z_\ell=0$ in relative holomorphic coordinates.

For a normal space $X$ and an integral Weil divisor $D$, $\OO_X(D)$
denotes the rank-one reflexive divisorial sheaf. We set
\[
 \omega_X^{[0]}=\OO_X,\qquad
 \omega_X^{[m]}=(\omega_X^{\otimes m})^{**}\quad(m\ge1).
\]
The symbol $K_X$ is additive notation for the canonical reflexive
sheaf. Thus, for an integral Weil divisor $D$ and $m\ge0$, we write
\[
 \OO_X(mK_X+D)
 :=\bigl(\omega_X^{[m]}\otimes\OO_X(D)\bigr)^{**}.
\]
Relative canonical notation has the same meaning. For a canonical
space, the $m$th plurigenus is $P_m(X)=h^0(X,\omega_X^{[m]})$.
A positive degree $m$ is \emph{admissible} for a rational boundary $B$
if $mB$ is integral; the associated adjoint sheaf is
$\OO_X(m(K_X+B))=\OO_X(mK_X+mB)$.

An adjoint $J$ is $\Q$-Cartier when a positive integral multiple
defines an invertible sheaf. Its Cartier multiples determine a
holomorphic $\Q$-line bundle. We write $\{J\}=c_1(J)$ for its
Bott--Chern class. We write $J_1\sim_{\Q}J_2$ when sufficiently
divisible Cartier multiples define isomorphic line bundles.
The adjoint $J$ is semiample if one such multiple is globally generated.
For a proper morphism $g:Y\to S$, it is relatively semiample if some
Cartier multiple $L=\OO_Y(qJ)$ has a surjective evaluation map
$g^*g_*L\to L$.

For a holomorphic map $u:T\to S$, write $Y_T=Y\times_S T$, with
projections $g_T:Y_T\to T$ and $p:Y_T\to Y$. Base change for a coherent
sheaf $F$ refers to the natural map
\[
 u^*R^ig_*F\longrightarrow R^i(g_T)_*(p^*F),
\]
where $p^*F$ is the ordinary sheaf pullback. The space $T$ may be
nonreduced. Identifications of these pullbacks with divisorial sheaves
on normal fibers will be established by adjunction.

Using compatible local meromorphic canonical frames, we write the
discrepancy formula on a log resolution as
\[
 K_T+B_T=\mu^*(K_X+B)+\sum_Q a(Q,X,B)Q,
\]
where $B_T$ is the strict transform and the sum is over exceptional
primes. At a nonexceptional boundary prime of coefficient $b$, set
$a=-b$. A pair is canonical, respectively terminal, when all exceptional
discrepancies are nonnegative, respectively positive. Klt means that
every discrepancy, including the nonexceptional ones, is greater than
$-1$. In particular, a canonical pair with coefficients less than one
is klt. Meromorphic computations are made on local Stein neighborhoods;
we use descent through proper modifications of normal spaces and finite
group quotients.

\subsection{Cohen--Macaulayness and restriction}

We recall the index-cover proof of Cohen--Macaulayness for divisorial
sheaves on klt spaces \cite[Corollary~5.25]{KM1998}. Applying it to a
normal Cartier fiber gives the restriction isomorphisms needed for
the infinitesimal thickenings.

\begin{lemma}\label{lem:CM}
Let $Z$ be a normal complex space with klt singularities, and let $D$
be an integral $\Q$-Cartier Weil divisor. Then $\OO_Z(D)$ is maximal
Cohen--Macaulay.
\end{lemma}

\begin{proof}
The assertion is local. Trivialize a Cartier multiple $rD$, and form
the normalized cyclic index cover. It is finite and étale in codimension
one, and its algebra has character pieces
\[
 \bigoplus_{a=0}^{r-1}\OO_Z(-aD).
\]
These pieces can be checked in codimension one and then reflexively
extended; this is the usual index-cover construction. The discrepancy
formula for a finite quasi-étale cover shows that the covering space is
klt. It has rational singularities and hence is Cohen--Macaulay.
A system of parameters at a point of $Z$ remains a system of parameters
at the finitely many points above it. The finite cover algebra is thus
a maximal Cohen--Macaulay module downstairs. Taking a character is a
direct summand operation in characteristic zero, so every displayed
summand is maximal Cohen--Macaulay. Local periodicity modulo $rD$
includes $\OO_Z(D)$.
\end{proof}

\begin{lemma}[Relative restrictions]\label{lem:restrictions}
Let $g:\cY\to\Delta$ be flat with normal central fiber $X$, and assume
that $\cY$ is normal and klt near $X$. Let $\cB$ be a rational horizontal
divisor, let $\cA=K_{\cY/\Delta}+\cB$ be $\Q$-Cartier, and suppose
adjunction gives $\cA|_X=A=K_X+B$. If $j\cB$ and $jB$ are integral,
then
\[
 \cF_j=\OO_{\cY}(j\cA)
\]
is flat over $\Delta$ and
\begin{equation}\label{eq:restriction}
 \cF_j/t\cF_j\cong\OO_X(jA).
\end{equation}
This isomorphism is induced by adjunction at the codimension-one
points of $X$. For $a\ge0$, put
\[
 R_a=\C[t]/(t^{a+1}),\qquad
 X_a=\cY\times_\Delta\Spec R_a,\qquad
 F_{j,a}=\cF_j/t^{a+1}\cF_j.
\]
These are flat $R_a$-sheaves, and for $a\ge1$ there is an exact sequence
\begin{equation}\label{eq:artin-basic}
 0\longrightarrow F_{j,a-1}\xrightarrow{\ t\ }F_{j,a}
 \longrightarrow F_{j,0}\longrightarrow0.
\end{equation}
For $\cB=0$ and $j=1$, these identifications give
$\OO_{\cY}(K_{\cY/\Delta})|_X\simeq\omega_X$.
\end{lemma}

\begin{proof}
A local meromorphic canonical frame represents $j\cA$ by an integral
$\Q$-Cartier Weil divisor. Lemma~\ref{lem:CM} therefore makes
$\cF_j$ maximal Cohen--Macaulay. The parameter
$t$ is a nonzerodivisor: the sheaf is torsion-free on the normal total
space and every component under consideration dominates the disk.
Torsion-freeness over the local discrete valuation ring of the base
implies flatness.

The quotient by $t$ is Cohen--Macaulay of full central dimension and
generically has rank one. It is therefore torsion-free and $S_2$, hence
reflexive on $X$. At a general codimension-one point of $X$, the fiber is smooth.
Flatness and smoothness of that fiber imply that $g$ is smooth near
the point. The relevant integral divisors are Cartier there, and the
usual adjunction and boundary restriction identify the quotient with
$\OO_X(jA)$. Since both sheaves are reflexive, this proves
\eqref{eq:restriction} globally. This argument also proves the
canonical-sheaf assertion directly.

Flatness supplies the displayed Artinian quotients and
\eqref{eq:artin-basic}. Equivalently, multiplication by $t$ embeds
$\cF_j/t^a\cF_j$ into $\cF_j/t^{a+1}\cF_j$ with quotient
$\cF_j/t\cF_j$. All identifications are compatible with restriction
and multiplication by the parameter.
\end{proof}

\begin{remark}\label{rem:nonCartier}
Equation~\eqref{eq:restriction} identifies the quotient itself with the
reflexive adjoint. This is the restriction used in the Artinian exact
sequences \eqref{eq:artin-basic}.
\end{remark}

\subsection{Birational comparison in every admissible degree}

\begin{lemma}\label{lem:birational}
Let $(Z_1,C_1)$ and $(Z_2,C_2)$ be normal canonical pairs with effective
boundaries, related by a
proper bimeromorphic correspondence. Assume that their boundaries
agree on the common nonexceptional prime divisors. Let $m\ge1$ with
$mC_1$ and $mC_2$ integral. On a common log resolution
$p_i:W\to Z_i$, let $C_W$ be the union of the strict transforms of
the two boundaries, with their prescribed coefficients. Then, in the
meromorphic pluricanonical identification,
\begin{equation}\label{eq:birational-sheaf}
 (p_i)_*\OO_W\bigl(m(K_W+C_W)\bigr)
 =\OO_{Z_i}\bigl(m(K_{Z_i}+C_i)\bigr),\qquad i=1,2.
\end{equation}
The statement holds on analytic total spaces over a disk, with relative
canonical divisors in place of the absolute ones.
\end{lemma}

\begin{proof}
For each $i$, the discrepancy formula is
\begin{equation}\label{eq:canonical-error}
 K_W+C_W=p_i^*(K_{Z_i}+C_i)+E_i,
 \qquad E_i\ge0,\quad E_i\text{ is }p_i\text{-exceptional}.
\end{equation}
Canonicity supplies the nonnegative exceptional discrepancies. A
boundary component present only on the other model adds a nonnegative
coefficient to the corresponding exceptional error. The nonexceptional
coefficients agree by hypothesis.

A local section downstairs has an effective $\Q$-Cartier zero divisor
relative to the integral adjoint $m(K_{Z_i}+C_i)$. Its pullback has
effective zero divisor after addition of $mE_i$. Hence the pulled-back
meromorphic $m$-canonical form lies in
$\OO_W(m(K_W+C_W))$. The resulting divisor is integral because
$W$ is smooth and $mC_W$ is integral.

Conversely, a section upstairs descends meromorphically and satisfies
the required pole bounds at every nonexceptional codimension-one point
of $Z_i$. The rank-one reflexive sheaf downstairs is determined by
these bounds, proving the reverse inclusion. Over a disk, tensoring
with a local generator of the base canonical line identifies the
relative calculation with the same absolute calculation.
\end{proof}

Over a disk, the total-space identity transports an extended section
along the correspondence. Its central value is determined on the common
dense smooth open set and hence everywhere by reflexivity.

\subsection{Persistence of an isomorphism locus}

The following elementary application of Nakayama's lemma will be
used in Sections~\ref{sec:kollar-family} and
\ref{sec:relative-reduction} to propagate central isomorphism loci
through deformed birational diagrams.

\begin{lemma}\label{lem:open-isomorphism}
Let $h:\mathcal Y\to\mathcal Z$ be a proper morphism over $\Delta$,
with $\mathcal Y$ flat over $\Delta$. If the induced morphism
$Y_0\to Z_0$ is an isomorphism over an open subset
$U_0\subset Z_0$, then $h$ is an isomorphism over a neighborhood
of $U_0$ in $\mathcal Z$.
\end{lemma}

\begin{proof}
Properness and upper semicontinuity of fiber dimension make $h$
finite near each point of $U_0$. Write $B\to A$ for its finite
local algebra there. The central hypothesis gives
$B/tB\simeq A/tA$, so Nakayama's lemma makes $B\to A$ surjective.
Let $I$ be its kernel. Since $A$ is flat over the base discrete
valuation ring, reduction of $0\to I\to B\to A\to0$ modulo $t$
is left exact. Thus $I/tI=0$, and Nakayama's lemma gives $I=0$.
The resulting neighborhoods cover $U_0$.
\end{proof}

\section{Central period injections on Kähler modifications}
\label{sec:central-periods}

The extension argument uses an injective period map on the central
fiber. We construct it on a compact Kähler modification and identify
it intrinsically on a smooth open subset. The Hodge-theoretic inputs
are logarithmic de Rham comparison
\cite[Proposition~3.1.8]{Deligne1971} and $E_1$-degeneration for compact
Kähler manifolds \cite[Theorem~3.2]{LRW2019}. For a compact Kähler
$n$-fold $Z$ and an SNC divisor $E$, they give the natural injection
\begin{equation}\label{eq:log-period-edge}
 H^1(Z,\omega_Z(E))\hookrightarrow H^{n+1}(Z\setminus E,\C).
\end{equation}
Naturality of this map under modifications will give the central
injections needed in Section~\ref{sec:endpoint}.

\subsection{The canonical degree}

\begin{lemma}\label{lem:period-one}
Let $X$ be a normal connected compact complex space with klt singularities.
Suppose there is a proper bimeromorphic morphism from a compact Kähler
manifold to $X$. Put $U=X_{\reg}$, $n=\dim X$, and
$\omega_X=\OO_X(K_X)$. Restriction, inclusion of the top holomorphic
de Rham term, and de Rham comparison define an injection
\begin{equation}\label{eq:period-one}
 \per_{1,0}:H^1(X,\omega_X)
 \longrightarrow H^1(U,\omega_U)
 \longrightarrow H^{n+1}(U,\C).
\end{equation}
\end{lemma}

\begin{proof}
Principalize the pullback ideal of $\Sing X$ on the given smooth Kähler
modification by projective blowups. This gives a proper modification
$p:Z\to X$, with $Z$ compact Kähler, such that
\[
 E=\bigl(p^{-1}(\Sing X)\bigr)_{\mathrm{red}}
\]
is an SNC divisor. If $X$ is smooth, take $E=\varnothing$. Put
$W=Z\setminus E$. The restriction $h:W\to U$ is a proper modification
between smooth spaces.

The sheaf calculation is
\begin{equation}\label{eq:period-one-sheaf}
 p_*\omega_Z(E)=\omega_X.
\end{equation}
For a local section $u$ of $\omega_X$, let $D_u$ be its effective
$\Q$-Cartier zero divisor. At a prime divisor $Q$ on $Z$, pullback gives
\[
 \ord_Q(p^*u)=v_Q(D_u)+a(Q,X,0).
\]
The order is nonnegative at nonexceptional divisors. At an exceptional
divisor it is an integer greater than $-1$, by the klt inequality, and
is again nonnegative. Thus $p^*u$ is holomorphic.

Conversely, let $\sigma$ be a section of $\omega_Z(E)$ over $p^{-1}(O)$,
where $O\subset X$ is open. Since $h$ is a proper modification of smooth
spaces, Hartogs extension gives $h_*\omega_W=\omega_U$. Hence $\sigma$
descends on $O\cap U$, and
$\omega_X=j_*\omega_U$, for $j:U\hookrightarrow X$, by reflexivity.
This proves \eqref{eq:period-one-sheaf}; the identification is pullback
of canonical forms.

The Leray edge map and \eqref{eq:log-period-edge} give the injection
\begin{equation}\label{eq:period-one-modification}
 H^1(X,\omega_X)\hookrightarrow
 H^1(Z,\Omega_Z^n(\log E))\hookrightarrow H^{n+1}(W,\C).
\end{equation}
The sheaf identity is ordinary pullback over $U$. Naturality of
restriction and de Rham comparison identifies
\eqref{eq:period-one-modification} with $h^*\circ\per_{1,0}$.
Hence $\per_{1,0}$ is injective.
\end{proof}

\subsection{Selected characters of a cyclic cover}

For higher adjoint powers, the cyclic-cover construction of
Esnault--Viehweg \cite[3.5, Corollary~3.11 and Lemma~3.16]{EV1992}
replaces ordinary canonical forms by a character eigensheaf. We first
construct the cover on a compact Kähler model, then identify this
eigensheaf by a divisorial calculation.

\begin{lemma}\label{lem:period-j}
Let $X$ be a normal connected compact complex space admitting a proper
bimeromorphic morphism from a compact Kähler manifold. Let $(X,B)$ be
klt, with $B$ effective rational and $A=K_X+B$ $\Q$-Cartier and
semiample. Choose a sufficiently divisible integer $k\ge2$ such that
$kB$ is integral and $kA$ is Cartier and globally generated,
and a general section $s_0\in H^0(X,\OO_X(kA))$ with zero divisor $G$,
such that
\[
 \left(X,B+\frac{k-1}{k}G\right)\quad\text{is klt}.
\]
Put
\[
 U=X_{\reg}\setminus\Supp(B+G).
\]
Let $q:V\to U$ be the finite étale $\boldsymbol\mu_k$-torsor of roots
of the nowhere vanishing $k$-canonical form $s_0|_U$. All its connected
components are included. Its tautological top form $\omega$ satisfies
$\omega^k=q^*s_0$. For each $2\le j\le k$
with $jB$ integral, put $F_j=\OO_X(jA)$ and define
\begin{equation}\label{eq:beta-j}
 \beta_j:F_j|_U\longrightarrow q_*\omega_V,\qquad
 \beta_j(u)=\frac{q^*u}{\omega^{j-1}}.
\end{equation}
Then restriction, $\beta_j$, and the top holomorphic de Rham term induce
an injection
\begin{equation}\label{eq:period-j}
 \per_{j,0}:H^1(X,F_j)\hookrightarrow H^{n+1}(V,\C),
 \qquad n=\dim X.
\end{equation}
\end{lemma}

\begin{proof}
We first justify the choice of $s_0$. Resolve $(X,B)$ after the given
compact Kähler modification, using projective blowups. The pullback
linear system $|kA|$ is generated. A general member is smooth, transverse
to the finitely many resolution strata, and contains no exceptional
center. Adding its pullback with coefficient $(k-1)/k<1$ to the crepant
boundary gives an SNC divisor whose coefficients are all less than one.
The asserted pair is therefore klt. The divisor $G$ is reduced and
shares no component with $B$. If the generated system is trivial, take
$G=0$. The same choice makes $(X,B+(j-1)G/k)$ klt for every $2\le j\le k$.

\emph{Step 1: a compact Kähler model of the cover.}
Put $S_X=X\setminus U$. Further projective principalization on the
chosen Kähler modification gives
\[
 \mu:T\longrightarrow X,\qquad
 S=(\mu^{-1}(S_X))_{\mathrm{red}}\text{ an SNC divisor},\qquad
 T_U=\mu^{-1}(U),
\]
where $T$ is smooth compact Kähler. The map $\mu_U:T_U\to U$ is a proper
modification of smooth manifolds.

The meromorphic $k$-canonical form $\mu^*s_0$ is holomorphic on $T_U$.
Choose an effective integral divisor $H$ supported on $S$ which clears
its poles, and put $L=\OO_T(K_T+H)$. The monic root hypersurface
of this section of $L^{\otimes k}$ is finite over $T$ and compact inside
$\PP(\OO_T\oplus L)$, disjoint from the infinity section. It is reduced:
the monic root algebra is locally free over $\OO_T$ and generically
separable, so any nilpotent would be torsion and must vanish.

Let $N_T$ be its normalization. There is a canonical identification
\begin{equation}\label{eq:root-normalization-over-open}
 (N_T)|_{T_U}=T_U\times_U V.
\end{equation}
To check it, write $s_0=f\tau^k$ on a coordinate open in $U$, with
$f$ a unit and $\tau$ a volume form. On an open in $T_U$ write
$\mu^*\tau=J\eta$, where $\eta$ is a volume form. The root equation is
$y^k=(f\circ\mu)J^k$. Its normalization is
$z^k=f\circ\mu$, via $y=Jz$: the latter cover is finite étale and
normal, and the finite map agrees with the root equation where $J\ne0$.
These local identifications glue to
\eqref{eq:root-normalization-over-open}.

Take an equivariant embedded resolution of the root hypersurface in
$\PP(\OO_T\oplus L)$ and resolve its inverse image of $S$. Its smooth
strict transform $Z$ is compact Kähler, since the ambient operations are
projective blowups. It factors through $N_T$. Let $p:Z\to X$ be the
composite and set
\[
 W=p^{-1}(U),\qquad E=(p^{-1}(S_X))_{\mathrm{red}}=Z\setminus W.
\]
We choose the resolution so that $E$ is SNC. Equation
\eqref{eq:root-normalization-over-open} gives a proper modification
\begin{equation}\label{eq:open-cover-modification}
 h:W\longrightarrow V,\qquad p|_W=q\circ h.
\end{equation}
It is the composite of the restricted resolution with the base change
of $\mu_U$, and is bimeromorphic on every component. Both $W$ and $V$
are smooth.

The cyclic action lifts to $Z$. Its tautological root, interpreted as
a meromorphic canonical form, gives a meromorphic top form on $Z$
restricting to $h^*\omega$ on $W$; denote it again by $\omega$.
Let $\chi$ be the character for which $\zeta^*\omega=\zeta\omega$.

\smallskip\noindent\emph{Step 2: the character eigensheaf.}
The required sheaf identity is
\begin{equation}\label{eq:selected-character-sheaf}
 \bigl(p_*\omega_Z(E)\bigr)_{\chi^{1-j}}=F_j.
\end{equation}
Over $U$ this identification is $h^*\beta_j$. Indeed,
$h_*\omega_W=\omega_V$ for the smooth proper modification $h$, by
the Hartogs argument used in Lemma~\ref{lem:period-one}. Consequently
\[
 (p_*\omega_Z(E))|_U=q_*h_*\omega_W=q_*\omega_V.
\]
The character summands of the degree-$k$ torsor are line bundles, and
\eqref{eq:beta-j} identifies its $\chi^{1-j}$ summand with $F_j|_U$.

For the forward inclusion in \eqref{eq:selected-character-sheaf}, let
$u$ be a local section of $F_j$, and let $D_u\ge0$ be its zero divisor
as a section of this reflexive adjoint sheaf. It is $\Q$-Cartier.
Put $\Theta_j=B+(j-1)G/k$. For a divisorial valuation $Q$ over $X$
and an extension $R$ upstairs with ramification index $e$, the order
of the meromorphic top form is
\begin{equation}\label{eq:selected-character-order}
 \ord_R\beta_j(u)
 =e\bigl(v_Q(D_u)+a(Q,X,\Theta_j)+1\bigr)-1.
\end{equation}
The discrepancy convention is that of Section~\ref{sec:sheaves}.
To verify the formula, realize $Q$ on a smooth model $\mu':T'\to X$
and put $a_Q=\operatorname{coeff}_Q(K_{T'}-\mu'^*A)$.
The pulled-back pluricanonical forms $u$ and $s_0$ have orders
\[
 A_Q=v_Q(D_u)+ja_Q,\qquad C_Q=v_Q(G)+ka_Q
\]
on $T'$. At a divisor of ramification index $e$, their further
pullbacks have orders $eA_Q+j(e-1)$ and $eC_Q+k(e-1)$.
Since $\omega^k=p^*s_0$, the order of $\omega$ is
$eC_Q/k+e-1$. Consequently
\[
\begin{aligned}
 \ord_R\beta_j(u)
 &=eA_Q+j(e-1)-(j-1)(eC_Q/k+e-1)\\
 &=e\bigl(v_Q(D_u)+a_Q-(j-1)v_Q(G)/k+1\bigr)-1,
\end{aligned}
\]
which is \eqref{eq:selected-character-order}.

To apply the calculation to an arbitrary prime of $Z$, use analytic
flattening of the proper generically finite map \cite{Hironaka1975}.
After a modification of the base, the strict transform is proper and
flat of relative dimension zero, hence finite. The prime then lies over
a divisor on the modified base. Normalization and resolution reduce to
the discrete valuation ring calculation above.

Since $(X,\Theta_j)$ is klt and $D_u$ is effective $\Q$-Cartier,
\eqref{eq:selected-character-order} is an integer greater than $-1$.
It is therefore nonnegative at every prime of $Z$. Thus $\beta_j(u)$
extends holomorphically across all divisors, and hence over $Z$.

For the reverse inclusion, let $\sigma$ be a local logarithmic top
form of character $\chi^{1-j}$ on $Z$. Over $W$, descent by $h$ and
the torsor identification express it uniquely as $\beta_j(u)$.
To extend $u$ meromorphically, choose a nonzero local meromorphic
canonical frame $\tau$ on $X$. The quotient
\begin{equation}\label{eq:meromorphic-character-ratio}
 \frac{\sigma\omega^{j-1}}{(p^*\tau)^j}
\end{equation}
is an invariant meromorphic function, hence descends meromorphically
to the normal quotient $Z/\boldsymbol\mu_k$. The induced map
$Z/\boldsymbol\mu_k\to X$ is proper and bimeromorphic: over a dense
open subset of $U$, the map $h$ is an equivariant isomorphism and
$V/\boldsymbol\mu_k=U$. Meromorphic descent through this modification,
followed by multiplication by $\tau^j$, gives $u$. For modification
descent, the proper image of the graph in $X\times\PP^1$ is analytic
and generically single-valued, hence defines the descended function.

At a generic boundary prime with coefficient $b\in(0,1)$, let $d$
be the integer order of $u$ relative to $F_j$. Logarithmic growth and
\eqref{eq:selected-character-order} imply
\[
 e(d+1-b)-1\ge-1,\qquad d\ge b-1>-1.
\]
Thus $d\ge0$. At a generic component of the reduced divisor $G$ the
same inequality gives
\[
 d\ge\frac{j-1}{k}-1>-1,
\]
again implying $d\ge0$. At every other prime divisor the cover is
generically étale and $E$ is absent, so no pole occurs. Modifications
of a normal space are isomorphisms at its generic codimension-one
points. Reflexivity of $F_j$ now proves the reverse inclusion in
\eqref{eq:selected-character-sheaf}.

\smallskip\noindent\emph{Step 3: the period map.}
By \eqref{eq:selected-character-sheaf}, the character projector
identifies $F_j$ with a direct summand of $p_*\omega_Z(E)$.
The Leray edge map and \eqref{eq:log-period-edge} therefore give
\begin{equation}\label{eq:period-j-modification}
 H^1(X,F_j)\hookrightarrow H^1(Z,\Omega_Z^n(\log E))
 \hookrightarrow H^{n+1}(W,\C).
\end{equation}
The sheaf map restricts on $W$ to $h^*\beta_j$. Naturality of the Leray
edge map and de Rham comparison identifies this composite with
\begin{equation}\label{eq:intrinsic-period-factorization}
 H^1(X,F_j)\xrightarrow{\per_{j,0}}H^{n+1}(V,\C)
 \xrightarrow{h^*}H^{n+1}(W,\C).
\end{equation}
Since this composite is injective, $\per_{j,0}$ is injective.
\end{proof}

\begin{remark}[The canonical character]\label{rem:central-period-scope}
The two period constructions use different open sets. For $j\ge2$,
the character bound along $G$ identifies the eigensheaf with $F_j$.
In degree one, the invariant logarithmic character allows a simple
pole along $G$, so Lemma~\ref{lem:period-one} instead uses
$U=X_{\reg}$. This also gives the degree-one injection for an arbitrary
klt central fiber with a compact Kähler modification.
\end{remark}

\section{Infinitesimal periods and analytic extension}
\label{sec:endpoint}

We apply Levine's infinitesimal extension method
\cite{Levine1983,Levine1985} to the central period maps of
Section~\ref{sec:central-periods}. First, a cohomological induction
converts the period injections into lifting over Artinian thickenings.
A finite-module argument then realizes these lifts analytically.

\subsection{An Artinian period criterion}

\begin{lemma}\label{lem:artin-period}
Let $g:\cY\to\Delta$ be proper and let $\cF$ be a coherent sheaf flat
over $\Delta$. For $a\ge0$, put
\[
 R_a=\C[t]/(t^{a+1}),\qquad
 X_a=\cY\times_\Delta\Spec R_a,\qquad
 F_a=\cF/t^{a+1}\cF,
\]
and write $X=X_0$. Fix $r\ge0$. Suppose there is a
complex vector space $H$ and, for $0\le a\le r$, $R_a$-linear maps
\[
 \Phi_a:H^1(X_a,F_a)\longrightarrow H\otimes_\C R_a
\]
compatible with truncation and with multiplication by $t$. If $\Phi_0$
is injective, then every $\Phi_a$ is injective and every restriction
\begin{equation}\label{eq:artin-surjection}
 H^0(X_a,F_a)\longrightarrow H^0(X,F_0)
\end{equation}
is surjective for $0\le a\le r$.
\end{lemma}

\begin{proof}
Flatness gives
\[
 0\longrightarrow F_{a-1}\xrightarrow{\ t\ }F_a
 \longrightarrow F_0\longrightarrow0.
\]
The statement at $a=0$ is the hypothesis together with the identity on
$H^0$. Suppose it is known at $a-1$. If a class
$b\in H^1(X_{a-1},F_{a-1})$ maps to zero under multiplication by $t$,
then $t\Phi_{a-1}(b)=0$. The map
\[
 H\otimes R_{a-1}\xrightarrow{\ t\ }H\otimes R_a
\]
is injective. Thus $\Phi_{a-1}(b)=0$, and $b=0$ by induction. The long
exact cohomology sequence gives the surjectivity in
\eqref{eq:artin-surjection}.

Next let $c\in H^1(X_a,F_a)$ satisfy $\Phi_a(c)=0$. Its central
restriction is zero by injectivity of $\Phi_0$. Exactness therefore
writes $c=tb$ for some $b\in H^1(X_{a-1},F_{a-1})$. Compatibility and
the same coefficient injection imply $\Phi_{a-1}(b)=0$. Hence $b=0$
and $c=0$. This completes both parts of the induction.
\end{proof}

For a line bundle on a smooth proper family, the passage from
finite-order lifting to constancy was proved in
\cite[Lemma~1.1]{Levine1983}. The next two lemmas give a
coherent-sheaf version of this passage.

\begin{lemma}[Analytic realization]\label{lem:analytic-lift}
With the notation and properness and flatness assumptions of
Lemma~\ref{lem:artin-period}, suppose that
$H^0(X_a,F_a)\to H^0(X,F_0)$ is surjective for every $a\ge0$.
Then every section of $F_0$ extends to a section of $\cF$ over the
inverse image of a neighborhood of $0$. The lifts at the separate
finite orders may be chosen independently.
\end{lemma}

\begin{proof}
Put
\[
 R=\C\{t\},\qquad M=(g_*\cF)_0,\qquad Q=(R^1g_*\cF)_0.
\]
Both are finite $R$-modules by Grauert's proper direct-image theorem
\cite{Grauert1960}. The direct-image sequence for multiplication by $t^a$ is
\begin{equation}\label{eq:finite-module}
 0\longrightarrow M/t^aM
 \longrightarrow H^0(X_{a-1},F_{a-1})
 \xrightarrow{\partial_a}Q[t^a]\longrightarrow0.
\end{equation}
Here the sheaf supported on the Artinian central thickening has direct
image stalk equal to the indicated cohomology group. The natural diagram
from multiplication by $t^a$ to multiplication by $t$ has multiplication
by $t^{a-1}$ on its left term and the identity on its middle term. Its
connecting maps give
\begin{equation}\label{eq:connecting-naturality}
 \partial_1\circ\operatorname{res}=t^{a-1}\partial_a.
\end{equation}

The torsion submodule of the finite module $Q$ is killed by $t^b$ for
some finite $b\ge0$. Choose $a>b$ and lift the prescribed central
section to $F_{a-1}$. Equations \eqref{eq:finite-module} and
\eqref{eq:connecting-naturality} show that its $\partial_1$ image is
zero. It therefore lies in the image of $M/tM$. A representative in
$M$ is, by the definition of a direct-image stalk, an analytic
section on the inverse image of a sufficiently small disk. This is the
desired extension. A finite basis of the central space gives a common
shrinking for all its sections.
\end{proof}

\begin{lemma}[Extension and base change]\label{lem:extension-constancy}
Let $g:\mathcal Y\to\Delta$ be proper and let $\mathcal F$ be coherent
and flat over $\Delta$. If every section of $\mathcal F|_{X_0}$ extends
near $X_0$, then $h^0(X_t,\mathcal F|_{X_t})$ is constant near $0$.
The direct image is locally free there and commutes with arbitrary
holomorphic base change.
\end{lemma}

\begin{proof}
The coherent sheaf $g_*\mathcal F$ is torsion-free on the disk, because
multiplication by a nonzero base germ is injective on $\mathcal F$.
It is therefore locally free. The direct-image sequence for multiplication
by $t$ injects its fiber at zero into $H^0(X_0,\mathcal F|_{X_0})$,
and the extension hypothesis makes this map surjective. Its rank thus
equals the central dimension. The corresponding injection at a nearby
point gives dimension at least this rank, while upper semicontinuity
gives dimension at most the central dimension. Equality follows near
zero. Grauert's theorem \cite{Grauert1960} gives base change.
\end{proof}

\subsection{Degree one}

\begin{proof}[Proof of Theorem~\ref{thm:degreeone}]
Let $\cF=\OO_{\cY}(K_{\cY/\Delta})$. Lemmas~\ref{lem:CM} and
\ref{lem:restrictions}, with the usual canonical adjunction, give
flatness and $F_0=\omega_X$. Put $U=X_{\reg}$ and let $U_a=X_a|_U$.
The morphism $U_a/\Spec R_a$ is smooth, and
\[
 F_a|_{U_a}=\omega_{U_a/R_a}.
\]
For every $a$, the relative holomorphic Poincaré lemma gives the
canonical quasi-isomorphism
\begin{equation}\label{eq:relative-poincare}
 \underline{R_a}\simeq\Omega^\bullet_{U_a/R_a},\qquad
 \HH^{n+1}(U_a,\Omega^\bullet_{U_a/R_a})
 =H^{n+1}(U,\C)\otimes R_a.
\end{equation}
Locally this is the holomorphic Poincaré lemma in relative coordinates
with coefficients in $R_a$; the canonical inclusions of constant scalars
glue to give the displayed quasi-isomorphism.

Restriction, the top term of the relative de Rham complex, and
\eqref{eq:relative-poincare} define
\[
 \Phi_a:H^1(X_a,F_a)\longrightarrow H^{n+1}(U,\C)\otimes R_a.
\]
They commute with multiplication by $t$ and truncation, and their central
map is injective by Lemma~\ref{lem:period-one}. Apply
Lemma~\ref{lem:artin-period} to every finite truncation and then
Lemma~\ref{lem:analytic-lift}.
\end{proof}

\subsection{Semiample adjoints}

The proof uses an auxiliary pluricanonical degree to control every
smaller admissible degree, following the argument that Levine credits
to N.~I.~Shepherd-Barron \cite[Corollary~1.8]{Levine1983}.
Lemma~\ref{lem:period-j} supplies the required character period maps
for reflexive adjoints.

\begin{proof}[Proof of Theorem~\ref{thm:endpoint}]
In degree $m=1$, admissibility and the coefficient bounds give
$\cB=B=0$, so Theorem~\ref{thm:degreeone} applies. A
zero-dimensional normal connected central fiber is a reduced point;
its flat deformation is smooth near that point, and the assertion
follows directly. Hence assume $m\ge2$ and $n=\dim X>0$.

\emph{Step 1: periods from an auxiliary lift.}
For $j=m$ and a sufficiently divisible $j=k$, set
\[
 \cF_j=\OO_{\cY}(j\cA),\qquad F_{j,a}=\cF_j/t^{a+1}\cF_j.
\]
Lemma~\ref{lem:restrictions} identifies $F_{j,0}$ with $\OO_X(jA)$
and gives the exact sequences of thickenings. Choose $k\ge m$ such
that $k\cA$ is Cartier near the compact central fiber, $k\cB$ is integral, and $kA$
is globally generated. A common Cartier multiple exists near the
whole central fiber by compactness, and a further multiple preserves
generation on $X$.

Choose $s_0$ and $G$ as in Lemma~\ref{lem:period-j}. Write
\[
 U=X_{\reg}\setminus\Supp(B+G),\qquad q:V\to U.
\]
Given a lift
\begin{equation}\label{eq:existing-lift}
 s_r\in H^0(X_r,F_{k,r})
\end{equation}
of $s_0$, use its compatible truncations $s_a$ for
$0\le a\le r$. On $U_a=X_a|_U$, the morphism is smooth and the
boundary is absent. Thus
\[
 F_{j,a}|_{U_a}=\omega_{U_a/R_a}^{\otimes j},
\]
and $s_a$ is a unit as a section of the $k$th canonical power. Its
canonical root cover
\[
 q_a:V_a\longrightarrow U_a,\qquad \omega_a^k=q_a^*s_a
\]
is finite étale and smooth over $R_a$. Its underlying topological
space is $V$, and the covers and tautological forms reduce compatibly
along this chosen finite chain.

By the same relative Poincaré calculation,
\[
 \HH^{n+1}(V_a,\Omega^\bullet_{V_a/R_a})
 =H^{n+1}(V,\C)\otimes R_a.
\]
The maps
\[
 \beta_{j,a}(u)=\frac{q_a^*u}{\omega_a^{j-1}}
\]
therefore define periods
\begin{equation}\label{eq:artin-periods}
 \per_{j,a}:H^1(X_a,F_{j,a})
 \longrightarrow H^{n+1}(V,\C)\otimes R_a,\qquad j=m,k.
\end{equation}
The forms $\omega_a$ are truncations of $\omega_r$, so these maps
commute with truncation and multiplication by $t$.

Lemma~\ref{lem:period-j} makes the central maps injective. Consequently
Lemma~\ref{lem:artin-period} gives, for this finite chain,
\begin{equation}\label{eq:period-and-sections}
 \per_{j,a}\text{ injective},\qquad
 H^0(X_a,F_{j,a})\twoheadrightarrow H^0(X,F_{j,0}),
 \qquad 0\le a\le r, j=m,k.
\end{equation}

Thus a lift of $s_0$ to order $r$ gives lifts of every section of
$\OO_X(mA)$ to that order. It remains to construct an auxiliary lift
for every $r$; Lemma~\ref{lem:analytic-lift} will then give the
analytic extension.

\smallskip\noindent\emph{Step 2: the auxiliary obstruction.}
Suppose that $s_r$ in \eqref{eq:existing-lift} has been constructed.
Let
\[
 \delta_{r+1}(s_0)\in H^1(X_r,F_{k,r})
\]
be the connecting class in
\begin{equation}\label{eq:aux-obstruction-sequence}
 0\longrightarrow F_{k,r}\xrightarrow{\ t\ }F_{k,r+1}
 \longrightarrow F_{k,0}\longrightarrow0.
\end{equation}
Its vanishing gives an order-$r+1$ lift of the central section $s_0$.
This is the inductive assertion: at each stage we require a lift of
$s_0$ to the indicated order.

Levine's differential calculation
\cite[Lemmas~1.2--1.5]{Levine1983} expresses this obstruction on $U_r$.
Choose a Stein coordinate cover of the smooth locus near $U$, and
transport sections to a common chart when taking Čech differences.
We use $(\check\delta a)_{ij}=a_i-a_j$. Write the relative transition
functions and their Jacobian determinants as
\[
 z_i=g_{ij}(z_j,t),\qquad
 J_{ij}=\det\left(\frac{\partial g_{ij}}{\partial z_j}\right).
\]
Differentiate the analytic transition functions before reducing
modulo $t^{r+1}$. In the target coordinates, put
\begin{equation}\label{eq:rho}
 \rho_{ij}=\sum_\alpha
 \frac{\partial g_{ij}^{\alpha}}{\partial t}
 \frac{\partial}{\partial z_i^\alpha}.
\end{equation}
Differentiating the coordinate cocycle identity shows that the resulting
relative vector fields form a Čech cocycle on $U_r$.

Write the chosen section on $U_r$ in canonical coordinates as
$s_r=h_i(dz_i)^k$. Choose arbitrary local lifts $\widetilde h_i$ of
its coefficients to order $r+1$. The gluing relation modulo $t^{r+1}$ reads
\[
 h_i(g_{ij}(z_j,t),t)J_{ij}^{k}=h_j(z_j,t).
\]
The chosen lifts therefore satisfy
\[
 \widetilde h_i(g_{ij}(z_j,t),t)J_{ij}^{k}-\widetilde h_j(z_j,t)
 =t^{r+1}\zeta_{ij}\pmod{t^{r+2}}.
\]
Its differentiated discrepancy, modulo $t^{r+1}$, leaves the cocycle
\begin{equation}\label{eq:levine-cocycle}
 D(s_r)_{ij}
 =\bigl(k\operatorname{div}(\rho_{ij})h_i
                  +\rho_{ij}(h_i)\bigr)(dz_i)^k.
\end{equation}
With this Čech convention, its class is
\begin{equation}\label{eq:levine-obstruction}
 [D(s_r)]=(r+1)\delta_{r+1}(s_0)|_{U_r}.
\end{equation}

Indeed, division of the difference of the chosen lifts by $t$ represents
the connecting class in \eqref{eq:aux-obstruction-sequence} by
$t^r\zeta_{ij}$. Differentiation represents $(r+1)$ times this class.
The local $\partial_t\widetilde h_i$ terms form a coboundary, leaving
\eqref{eq:levine-cocycle}.

\smallskip\noindent\emph{Step 3: vanishing of the period.}
Finite étaleness lifts $\rho_{ij}$ to relative vector fields on $V_r$.
Their contractions with the global tautological form,
\[
 \mu_{ij}=\iota_{\rho_{ij}}\omega_r,
\]
form a Čech cocycle of relative $(n-1)$-forms. Since
$q_r^*s_r=\omega_r^k$, Lie differentiation and Cartan's formula give
\begin{equation}\label{eq:exact-obstruction}
\begin{aligned}
 \beta_{k,r}\bigl(D(s_r)_{ij}\bigr)
 &=\frac{\mathcal L_{\rho_{ij}}(\omega_r^k)}{\omega_r^{k-1}}\\
 &=k\,\mathcal L_{\rho_{ij}}\omega_r
 =k\,\dd\mu_{ij}.
\end{aligned}
\end{equation}
The last expression is a boundary in the total Čech--de Rham complex,
because the Čech differential of $\mu$ is zero. Hence
\eqref{eq:levine-obstruction} has zero period. Injectivity of
$\per_{k,r}$ in \eqref{eq:period-and-sections}, and characteristic
zero, imply
\begin{equation}\label{eq:aux-obstruction-zero}
 \delta_{r+1}(s_0)=0.
\end{equation}

Starting with $s_0$, induction gives an auxiliary lift at every finite
order. Step~1 lifts every section of $\OO_X(mA)$ to each such order,
and Lemma~\ref{lem:analytic-lift} realizes its analytic extension.
\end{proof}

\begin{remark}\label{rem:scope-endpoint}
The Kähler modification provides the central period injections.
The relative obstruction calculation takes place on the smooth
finite-order covers of $U_r$; injectivity then detects the global
obstruction on $X_r$.
\end{remark}

\subsection{Canonical central fibers}

\begin{corollary}\label{cor:canonical-endpoint}
Let $g:\mathcal Y\to\Delta$ be proper and flat. Assume that the total
space is normal and canonical near the connected central fiber $X$,
and that $X$ is normal and canonical. Suppose $X$ admits a compact
Kähler modification and $K_X$ is semiample. Then for every fixed
$m\ge1$, every section of $\omega_X^{[m]}$ extends near $X$.
After shrinking so that the nearby fibers are normal and canonical,
\[
 t\longmapsto h^0(X_t,\omega_{X_t}^{[m]})
\]
is constant, and the corresponding direct image is locally free and
satisfies holomorphic base change.
\end{corollary}

\begin{proof}
Canonical singularities are klt, and the canonical divisor is
$\Q$-Cartier. The usual adjunction for a normal Cartier fiber identifies
$K_{\mathcal Y/\Delta}|_X$ with $K_X$. Thus Theorem~\ref{thm:endpoint},
with boundary zero, extends every reflexive canonical power.
The analytic deformation theorem for canonical singularities
\cite[Lemma~10 and Remark~10.3]{Kollar2021} permits a shrinking on which
all fibers are normal and canonical. Lemma~\ref{lem:restrictions},
applied at each such fiber, identifies the restriction of
$\OO_{\mathcal Y}(mK_{\mathcal Y/\Delta})$ with
$\omega_{X_t}^{[m]}$. Lemma~\ref{lem:extension-constancy} now gives
the remaining conclusions for the fixed degree $m$.
\end{proof}

In degree one, Theorem~\ref{thm:degreeone} gives the same extension
conclusion without semiampleness.


\section{Uniform neighborhoods and graded algebras}
\label{sec:uniform}

A generated multiple groups all adjoint powers into finitely many
residue classes. Pushing these classes to relative projective space
gives one neighborhood for all degrees and controls the multiplication
maps. This is the finite-generation
principle used in \cite[Lemma~11 and paragraph~13]{Kollar2021}.
We first treat global sections, then graded algebras and higher
cohomology.

\begin{lemma}[Uniformity from a generated multiple]
\label{lem:uniform}
Let $g:\mathcal Y\to\Delta$ be proper and flat. Suppose that coherent
sheaves $\mathcal F_m$, $m\ge0$, are flat over $\Delta$, with
$\mathcal F_0=\OO_{\mathcal Y}$. Suppose that, for an integer $k>0$,
$L=\mathcal F_k$ is invertible and there are identifications
\[
 \mathcal F_{ak+r}=\mathcal F_r\otimes L^a,
 \qquad a\ge0,\quad 0\le r<k,
\]
compatible with fiber restriction. Assume that $L_0$ is globally
generated and that every central section of every $\mathcal F_m$,
$m\ge1$, extends to a germ over the disk. Then there is a disk about $0$
on which all the dimensions
\[
 h^0(Y_t,\mathcal F_m|_{Y_t}),\qquad m\ge1,
\]
are constant. On that disk $g_*\mathcal F_m$ is locally free and commutes with
holomorphic base change for every $m\ge1$.
\end{lemma}

\begin{proof}
\emph{Step 1: a generated multiple.}
Extend a finite basis of $H^0(Y_0,L_0)$ and shrink once. The cokernel
of the resulting evaluation $\OO_{\mathcal Y}^{N+1}\to L$ vanishes
at every central point by Nakayama's lemma. Its support is closed;
properness of $g$ allows a further shrinking on which it is empty.
Consequently these sections define a proper morphism
\[
 \phi:\mathcal Y\longrightarrow P:=\PP^N\times\Delta,
 \qquad L=\phi^*\OO_P(1).
\]
Properness follows from the closed graph and the properness of $g$.

\smallskip\noindent\emph{Step 2: base change for the residue sheaves.}
For the finite set of residues $0\le r<k$, put
\[
 E_r=\phi_*\mathcal F_r,\qquad Q_r=R^1\phi_*\mathcal F_r.
\]
They are coherent. Each $E_r$ is torsion-free over the disk, hence flat.
At the origin, the long exact sequence for multiplication by $t$ gives
\begin{equation}\label{eq:residue-defect}
 0\longrightarrow E_r/tE_r
 \longrightarrow (\phi_0)_*(\mathcal F_r|_{Y_0})
 \longrightarrow T_r\longrightarrow0,
 \qquad T_r=Q_r[t].
\end{equation}
We prove $T_r=0$ by extending sections in one sufficiently large
degree $ak+r$ for each residue $r$.

Fix $r$ and choose one sufficiently large integer $a$ such that
\[
 H^1(\PP^N,(E_r/tE_r)(a))=0,
 \qquad T_r(a)\ \text{is globally generated},\qquad ak+r>0.
\]
Serre vanishing and generation, applied through GAGA
\cite{Serre1956}, give such an $a$. Every section of
$\mathcal F_{ak+r}|_{Y_0}$ extends by hypothesis. The projection
formula makes its extension a section of $E_r(a)$ near the central
projective space. Therefore
\[
 H^0(\PP^N,(E_r/tE_r)(a))
 \longrightarrow
 H^0(\PP^N,(\phi_0)_*(\mathcal F_r|_{Y_0})(a))
\]
is surjective. Taking cohomology in \eqref{eq:residue-defect} and
using the chosen $H^1$ vanishing gives $H^0(T_r(a))=0$.
Global generation forces $T_r=0$.

The equality $Q_r[t]=0$ says that $Q_r$ is flat over the disk at
every point of $P_0$. Relative flatness of a coherent analytic
sheaf is open \cite{Frisch1967}. The projection $P\to\Delta$ is proper, so we may
shrink once to make all the finitely many $Q_r$ flat over the whole
new disk. At any point $c$ of this disk, the multiplication-by-$t-c$
sequence now gives
\begin{equation}\label{eq:residue-base-change}
 E_r|_{P_c}\simeq(\phi_c)_*(\mathcal F_r|_{Y_c}).
\end{equation}
\smallskip\noindent\emph{Step 3: uniform regularity.}
Choose one nonnegative integer $a_0$ for which all the central
sheaves $E_r|_{P_0}$ are $a_0$-regular, that is,
\[
 H^i(\PP^N,E_r|_{P_0}(a_0-i))=0
 \quad(1\le i\le N,\ 0\le r<k).
\]
These are finitely many conditions. Upper semicontinuity and one
further shrinking make them hold at every nearby $c$.
Castelnuovo--Mumford regularity \cite[Lecture~14]{Mumford1966} then gives
\begin{equation}\label{eq:uniform-serre}
 H^i(\PP^N,E_r|_{P_c}(a))=0
 \quad(i>0,\ a\ge a_0,\ 0\le r<k)
\end{equation}
simultaneously on this disk. For $N=0$ there is no higher
cohomology and the same conclusion is immediate.

Flatness of $E_r$ makes the Euler characteristic of each fixed
fiberwise twist constant on the connected disk. Together
with \eqref{eq:uniform-serre}, this makes
$h^0(\PP^N,E_r|_{P_c}(a))$ independent of $c$ for every $a\ge a_0$.
By \eqref{eq:residue-base-change} and the projection formula these
are exactly the dimensions in degrees $ak+r$. Only finitely many
positive degrees remain. Lemma~\ref{lem:extension-constancy} handles
them, followed by a finite intersection of neighborhoods.
Grauert's theorem gives the asserted local freeness and base change
on the already fixed disk, for every degree.
\end{proof}

\begin{corollary}[Semiample canonical central fibers]
\label{cor:uniform-endpoint}
Let $g:\mathcal Y\to\Delta$ be proper and flat. Suppose that the
connected central fiber $Y_0$ is normal and canonical, admits a
compact Kähler modification, and has semiample canonical divisor.
After shrinking once, the total space and all fibers are canonical,
$K_{\mathcal Y/\Delta}$ is relatively semiample, and
\[
 h^0(Y_t,\omega_{Y_t}^{[m]})=h^0(Y_0,\omega_{Y_0}^{[m]})
 \qquad(t\in\Delta,\ m\ge1).
\]
For every $m\ge1$, the sheaf
$g_*\OO_{\mathcal Y}(mK_{\mathcal Y/\Delta})$ is locally free and
commutes with holomorphic base change on this disk. Every central
section extends over it.
\end{corollary}

\begin{proof}
By \cite[Lemma~10 and Remark~10.3]{Kollar2021}, shrink so that the
total space and every fiber are normal canonical. The nearby fibers
are connected by upper semicontinuity of $h^0(\OO)$.
Take a common relative Cartier index and a multiple $k$ for which
$\omega_{Y_0}^{[k]}$ is generated. With
\[
 \mathcal F_m=\OO_{\mathcal Y}(mK_{\mathcal Y/\Delta}),
 \qquad L=\OO_{\mathcal Y}(kK_{\mathcal Y/\Delta}),
\]
the restrictions and flatness follow from
Lemma~\ref{lem:restrictions}. Tensoring the reflexive sheaves
$\mathcal F_r$ with $L^a$ gives the residue identities. Central extension
in every degree is Corollary~\ref{cor:canonical-endpoint}.
Apply Lemma~\ref{lem:uniform}. Its first step makes
$L=\OO_{\mathcal Y}(kK_{\mathcal Y/\Delta})$ relatively generated,
which proves relative semiampleness. Finally, the locally free
direct image on the Stein disk extends every vector in its central
fiber to a section over the disk.
\end{proof}


\begin{corollary}[Uniform logarithmic endpoint]
\label{cor:uniform-log-endpoint}
Under the assumptions of Theorem~\ref{thm:endpoint}, let
$d>0$ be the least integer for which both $d\mathcal B$ and $dB$ are integral.
After shrinking, write $B_t$ for the boundary given by
restriction and adjunction on the normal fiber $Y_t$.
Then there is a disk about $0$ on which, for every fiber,
\[
 h^0(Y_t,\OO_{Y_t}(m(K_{Y_t}+B_t)))
 =h^0(X,\OO_X(m(K_X+B)))\qquad(m\in d\mathbb Z_{>0}).
\]
The sheaves $g_*\OO_{\mathcal Y}(m\mathcal A)$ are locally free
and commute with holomorphic base change for all such $m$.
Moreover, a positive Cartier
multiple of $K_{\mathcal Y/\Delta}+\mathcal B$ is generated
over the disk.
\end{corollary}

\begin{proof}
Normality of the fibers is open, so we first shrink around the
normal central fiber. Since $\mathcal Y$ is klt,
$K_{\mathcal Y/\Delta}$ is $\mathbb Q$-Cartier; hence so is
$\mathcal B=\mathcal A-K_{\mathcal Y/\Delta}$.
Its restriction defines $B_t$ and the usual adjunction gives
$\mathcal A|_{Y_t}=K_{Y_t}+B_t$.
At a general codimension-one point of a normal fiber the
morphism is smooth, so $dB_t$ is integral.
Lemma~\ref{lem:restrictions} therefore applies to every fiber.

Use $\mathcal F_j=\OO_{\mathcal Y}(jd\mathcal A)$.
Theorem~\ref{thm:endpoint} extends every central section of
every $\mathcal F_j$. Choose $k$ such that $kd\mathcal A$
is Cartier and $kdA$ is globally generated. The required
residue identities follow by tensoring with the invertible
$\mathcal F_k$. Lemma~\ref{lem:uniform} and its generation
argument prove the assertions.
\end{proof}

\subsection{The graded algebra}

\begin{proposition}[The graded algebra at a semiample endpoint]
\label{prop:endpoint-graded-algebra}
In the setting of Lemma~\ref{lem:uniform}, assume in addition that
$\mathcal Y$ is normal, the central fiber is normal and connected,
and the sheaves $\mathcal F_m$ form a commutative graded
$\OO_{\mathcal Y}$-algebra with degree-zero part $\OO_{\mathcal Y}$.
Assume that the residue identifications in that lemma are induced by
its multiplication. After shrinking to one disk $U$, the algebra
\[
 \mathcal R=\bigoplus_{m\ge0}g_*\mathcal F_m
\]
is finitely generated and flat over $\OO_U$, and every graded piece
is locally free. In the base-change notation of
Section~\ref{sec:sheaves}, the natural map for every holomorphic
$u:T\to U$ is an isomorphism of graded algebras:
\begin{equation}\label{eq:graded-holomorphic-base-change}
 u^*\mathcal R\simeq
 \bigoplus_{m\ge0}(g_T)_*(p^*\mathcal F_m).
\end{equation}
The relative analytic space
\[
 \mathcal C=\operatorname{Proj}_U\mathcal R
\]
is normal, projective and flat over $U$, and its formation commutes
with these base changes.
\end{proposition}

\begin{proof}
Use the disk and notation of Lemma~\ref{lem:uniform}. In particular,
$L=\mathcal F_k$ is relatively generated,
$\phi:\mathcal Y\to P=\PP^N\times U$ satisfies
$L=\phi^*\OO_P(1)$, and
$E_r=\phi_*\mathcal F_r$ for $0\le r<k$.
After shrinking, $g_*\OO_{\mathcal Y}=\OO_U$: upper
semicontinuity gives $h^0(Y_t,\OO_{Y_t})=1$, and Grauert's theorem
identifies the unit map with an isomorphism.

\emph{Finite generation and base change.}
The uniform Castelnuovo--Mumford regularity used in that lemma also
gives, for every $t\in U$, $r<k$ and $a\ge a_0$, the surjection
\[
 H^0(\PP^N,\OO(1))\otimes H^0(P_t,E_r|_{P_t}(a))
 \longrightarrow H^0(P_t,E_r|_{P_t}(a+1));
\]
see \cite[Lecture~14]{Mumford1966}. The restriction isomorphisms
\eqref{eq:residue-base-change}, the projection formula, and
Lemma~\ref{lem:uniform} identify this with the fiber of
\[
 \OO_U^{N+1}\otimes\mathcal R_{ak+r}
 \longrightarrow\mathcal R_{(a+1)k+r}.
\]
Its cokernel vanishes by Nakayama's lemma. Thus the finitely many
coherent sheaves $\mathcal R_{ak+r}$ with $0\le a\le a_0$ and
$0\le r<k$ generate $\mathcal R$ as a module over
$\operatorname{Sym}_{\OO_U}(\OO_U^{N+1})$, whose generators have
degree $k$. This proves finite generation after one further
shrinking. The algebra is flat over $U$, being a direct sum of
locally free graded pieces. Degreewise holomorphic base change is
Lemma~\ref{lem:uniform}; its naturality with respect to multiplication
proves \eqref{eq:graded-holomorphic-base-change}.

\smallskip\noindent\emph{The projective model.}
Take the Stein factorization
\[
 \mathcal Y\xrightarrow{h}\mathcal C'
 \xrightarrow{v}P
\]
of $\phi$. The normality of $\mathcal Y$ implies the normality of
$\mathcal C'$. Set $A=v^*\OO_P(1)$, which is ample over $U$
because $v$ is finite. Since $h_*\OO_{\mathcal Y}=\OO_{\mathcal C'}$
and $L=h^*A$, the $k$th Veronese algebra is
\[
 \mathcal R^{(k)}
 =\bigoplus_{a\ge0}g_*L^a
 =\bigoplus_{a\ge0}(g_{\mathcal C'})_*A^a.
\]
The relative Proj of the section algebra of an ample line bundle
recovers its space. Relative Proj is unchanged by a Veronese
subalgebra, so $\mathcal C=\mathcal C'$. This also proves projectivity.
After shrinking, the normal total space $\mathcal Y$ is irreducible
and dominates $U$; the same holds for $\mathcal C'$. Its local rings
are therefore torsion-free over the local discrete valuation rings
of $U$, and $\mathcal C\to U$ is flat.

Finally, after taking a sufficiently divisible Veronese, a finite
homogeneous presentation realizes $\mathcal C$ as a closed analytic
subspace of a projective bundle over $U$. Pulling back its homogeneous
equations proves that relative analytic Proj commutes with arbitrary
holomorphic base change. Together with
\eqref{eq:graded-holomorphic-base-change}, this identifies every
base-changed projective model.
\end{proof}

\subsection{Higher cohomological degrees}

\begin{lemma}[Uniformity in all cohomological degrees]
\label{lem:uniform-higher}
Let $g:\mathcal Y\to\Delta$ be proper and flat, and let
$\mathcal F_m$, $m\ge0$, be coherent sheaves flat over $\Delta$,
with $\mathcal F_0=\OO_{\mathcal Y}$.
Suppose that $L=\mathcal F_k$ is invertible for some $k>0$,
that $L_0$ is globally generated, and that
\[
 \mathcal F_{ak+r}\simeq\mathcal F_r\otimes L^a
 \qquad(a\ge0,\ 0\le r<k),
\]
compatibly with fiber restriction. Assume that, for every fixed
$m\ge1$, all the functions
$t\mapsto h^i(Y_t,\mathcal F_m|_{Y_t})$ are constant near $0$.
Then there is a disk about $0$ on which all these dimensions are
constant simultaneously for $m\ge1$ and $i\ge0$. Every
$R^ig_*\mathcal F_m$, with $m\ge1$, is locally free there and
commutes with holomorphic base change.
\end{lemma}

\begin{proof}
Grauert's theorem in degree $k$ extends a basis of $H^0(Y_0,L_0)$.
As in Lemma~\ref{lem:uniform}, its evaluation generates $L$ after
shrinking and defines
\[
 \phi:\mathcal Y\to P=\PP^N\times\Delta,
 \qquad L=\phi^*\OO_P(1),\qquad g=\pi\phi.
\]
For $0\le r<k$, set $Q_{r,j}=R^j\phi_*\mathcal F_r$. Proper
coherence gives a finite collection of coherent sheaves: the terms
with $j>\dim Y_0$ vanish.

\emph{Step 1: flatness of the higher direct images.}
Choose one integer $a>0$, and shrink the disk, so that
\begin{equation}\label{eq:higher-large-twist}
 R^p\pi_*Q_{r,j}(a)=0\quad(p>0),
 \qquad Q_{r,j}[t](a)\text{ is globally generated on }P_0
\end{equation}
for every $r,j$. To make this choice, consider the $t$-power
torsion subsheaf of each $Q_{r,j}$ near the compact fiber $P_0$. Coherence
and compactness give a common annihilating power of $t$, so this
torsion has a finite filtration with coherent quotients on $P_0$.
The torsion-free quotient is flat at $P_0$ and, after shrinking,
flat over the disk \cite{Frisch1967}. Apply Serre vanishing to
the filtration quotients and to the central fiber of the flat
quotient. Upper semicontinuity and Grauert's theorem give the
first part of \eqref{eq:higher-large-twist} for one sufficiently
large $a$. Serre generation for the finitely many $Q_{r,j}[t]$
gives the second part.

The projection formula and Leray now identify
\[
 \pi_*Q_{r,j}(a)\simeq R^jg_*\mathcal F_{ak+r}.
\]
The right side is locally free at the origin by the fixed-degree
hypothesis and Grauert's theorem. The $t$-torsion of the left side
is $H^0(P_0,Q_{r,j}[t](a))$. It is zero; generation in
\eqref{eq:higher-large-twist} therefore gives $Q_{r,j}[t]=0$.
Thus all $Q_{r,j}$ are flat at $P_0$. Openness of flatness and
properness permit one shrinking on which they are all flat.

\smallskip\noindent\emph{Step 2: fiberwise base change.}
For each $c$ in this disk, the direct-image sequence for
multiplication by $t-c$ gives
\[
 0\longrightarrow Q_{r,j}/(t-c)Q_{r,j}
 \longrightarrow R^j(\phi_c)_*(\mathcal F_r|_{Y_c})
 \longrightarrow Q_{r,j+1}[t-c]\longrightarrow0.
\]
The last term vanishes by flatness. Hence
\begin{equation}\label{eq:higher-residue-base-change}
 Q_{r,j}|_{P_c}\simeq
 R^j(\phi_c)_*(\mathcal F_r|_{Y_c})
\end{equation}
for every $r,j$ and $c$.

\smallskip\noindent\emph{Step 3: uniform regularity.}
Choose a common regularity bound $a_0$ for the finitely many
central sheaves $Q_{r,j}|_{P_0}$. Their flatness and upper
semicontinuity make the same bound valid for every nearby fiber.
Castelnuovo--Mumford regularity \cite[Lecture~14]{Mumford1966}
gives
\[
 H^p(P_c,Q_{r,j}|_{P_c}(a))=0
 \qquad(p>0,\ a\ge a_0).
\]
Consequently, fiberwise Leray and
\eqref{eq:higher-residue-base-change} identify
\[
 H^j(Y_c,\mathcal F_{ak+r}|_{Y_c})
 \simeq H^0(P_c,Q_{r,j}|_{P_c}(a))
 \qquad(a\ge a_0).
\]
The dimensions on the right equal Euler characteristics, which
are constant because $Q_{r,j}$ is flat. This treats every large
degree on one disk. Only finitely many positive degrees remain;
their fixed-degree neighborhoods have a common intersection.
Grauert's theorem proves local freeness and base change there.
\end{proof}

\section{Canonical negative steps}
\label{sec:canonical-steps}

The central program must start on the given canonical fiber. We use
the steps of \cite[Definition~7]{Kollar2021}: diagrams
\[
 X\xrightarrow f Z\xleftarrow q X^+
\]
of projective bimeromorphic contractions such that $-K_X$ is
$f$-ample, $q$ is small, and $K_{X^+}$ is $q$-ample. We call such a
diagram a \emph{canonical step}; it includes the mixed steps of
\cite[Definition~7]{Kollar2021}. The negative contraction may
contract several divisors. Throughout this section the spaces are
normal and irreducible.

\subsection{The divisorial canonical algebra}

We first recall the divisorial-algebra construction
\cite[3.52 and Lemma~6.2]{KM1998} over a neighborhood of a prescribed
compact subset of the target. It supplies the positive side of a canonical step; compare
\cite[Definition~7]{Kollar2021}.

\begin{lemma}\label{lem:canonical-algebra}
Let $Z$ be a normal complex analytic space whose divisorial canonical
algebra
\[
 \mathcal R_Z=\bigoplus_{m\ge0}\omega_Z^{[m]}
\]
is locally finitely generated, and let $K\subset Z$ be compact.
There are an open neighborhood $U$ of $K$ and an integer $r>0$ such that,
for
\[
 q:Y=\operatorname{Proj}_U(\mathcal R_Z|_U)\longrightarrow U,
\]
$Y$ is normal, $q$ is small, projective and bimeromorphic, and
$\omega_Y^{[r]}$ is invertible and $q$-ample. In particular, $K_Y$
is $\Q$-Cartier and $q$-ample. If $Z$ is compact, one may take $U=Z$.
\end{lemma}

\begin{proof}
Local finite generation supplies, around each point of $K$, an open
set on which a sufficiently divisible Veronese is generated in degree
one. Choose finitely many such sets $V_i$ covering $K$, and let $r$
be a common multiple of their Veronese degrees. On $U=\bigcup_iV_i$,
the tautological sheaf $L=\OO_Y(r)$ is invertible and relatively ample.
Thus $q$ is projective.

Normality can be checked over a connected Stein open set $V\subset U$.
Cartan's theorem~A gives a section of $\omega_V$ nonzero at a chosen smooth point;
it supplies a local meromorphic canonical trivialization.
Write its canonical Weil divisor as $\sum_P k_PP$.
In the meromorphic function field with an added variable $T$, the
divisorial algebra is the intersection of the polynomial ring with
the valuation rings defined by the Gauss valuations
\[
 w_P\left(\sum_m a_mT^m\right)
   =\min_m\bigl(v_P(a_m)+mk_P\bigr).
\]
The inequalities $w_P\ge0$ express precisely the codimension-one
regularity of each homogeneous coefficient. This intersection is
integrally closed. The degree-zero rings after homogeneous localization
are consequently normal, so the relative analytic Proj is normal.
The same construction in the excellent analytic local rings gives
normal analytic Proj charts. Over the open set where the canonical sheaf is
invertible, this Proj is the original space. Thus it is bimeromorphic.

The relative-Proj comparison holds locally for all sufficiently large
degrees. A finite covering of $K$, followed by a shrinking of $U$,
gives one bound $n_0$ such that
\begin{equation}\label{eq:canonical-proj-push}
 q_*L^n=\omega_U^{[rn]}
 \qquad(n\ge n_0).
\end{equation}
Suppose that $E$ is a $q$-exceptional prime divisor. Viewing all the
following sheaves in the same meromorphic rank-one sheaf, normality
and reflexivity give
\[
 q_*L^n\subseteq q_*(L^n\otimes\OO_Y(E))
 \subseteq(q_*L^n)^{**}
 =\omega_U^{[rn]}=q_*L^n
 \qquad(n\ge n_0).
\]
The second inclusion holds because the sheaves agree off $q(E)$,
which has codimension at least two. Hence all inclusions are equalities.
Relative Serre generation makes $L^n\otimes\OO_Y(E)$ generated over
a neighborhood of any fixed point of $q(E)$ for sufficiently large $n$.
Its evaluation map, however, now factors through
the proper subsheaf $L^n\subsetneq L^n\otimes\OO_Y(E)$ at the
generic point of $E$. This contradiction proves smallness.

The equality $L=\omega_Y^{[r]}$ holds on the isomorphism locus of
$q$. Both sheaves are reflexive, and the complement has codimension
at least two, so it holds everywhere. Thus $K_Y$ is $\Q$-Cartier
and $q$-ample, with the single index $r$ on $Y$.
\end{proof}

\begin{lemma}\label{lem:canonical-replacement}
Let $f:X\to Z$ be a projective bimeromorphic contraction of normal
analytic spaces. Suppose $X$ is canonical and $-K_X$ is $f$-ample.
Fix a compact subset $K\subset Z$. There is an open neighborhood
$U$ of $K$ on which the following assertions hold. Replace $Z$ by $U$
and $X$ by $f^{-1}(U)$, retaining the notation $f:X\to Z$ for the
restricted morphism. Then
\begin{equation}\label{eq:negative-canonical-algebra}
 f_*\omega_X^{[m]}=\omega_Z^{[m]}\qquad(m\ge0).
\end{equation}
The relative Proj of this algebra completes $f$ to a canonical step
$X\dashrightarrow X^+$. The space $X^+$
is canonical. On a common projective resolution $p:W\to X$,
$a:W\to X^+$, there is a $\Q$-linear equivalence of holomorphic
$\Q$-line bundles
\begin{equation}\label{eq:mixed-canonical-comparison}
 p^*K_X\sim_{\Q}a^*K_{X^+}+E_f,
 \qquad E_f\ge0,\qquad a_*E_f=0.
\end{equation}
If $f$ is nontrivial, $E_f\ne0$; its coefficient is positive at the
strict transform of every $f$-exceptional prime of $X$.
The space $Z$ has rational singularities. If the original target is
compact Kähler, no restriction is necessary, and $X^+$ is compact Kähler.
\end{lemma}

\begin{proof}
For \eqref{eq:negative-canonical-algebra}, work over a connected
Stein neighborhood $V\subset Z$. Choose a meromorphic canonical frame on $V$ as in the preceding
proof, and pull it back to $f^{-1}(V)$. A reflexively regular
$m$-canonical form on $Z$ defines a $\Q$-Cartier divisor
$D_X=\operatorname{div}(\varphi)+mK_X$ on $X$, with
$f_*D_X\ge0$ and $-D_X$ relatively nef. Exceptional negativity
gives $D_X\ge0$. Conversely, a form regular on $X$ is regular at
the generic points of the prime divisors of $Z$, and hence belongs
to its reflexive canonical power. This proves \eqref{eq:negative-canonical-algebra}.

Fujino's analytic form of canonical-algebra finite generation
\cite[Theorem~1.18]{Fujino2022}, extending
\cite[Corollary~1.1.2]{BCHM2010}, applies to $(X,0)$ and $f$.
Apply Lemma~\ref{lem:canonical-algebra} to the prescribed compact
subset $K$. After restriction to the resulting neighborhood and with
the retained notation, it gives a small projective $q:X^+\to Z$
with a single Cartier multiple of $K_{X^+}$ that is $q$-ample.
On the common resolution, subtract the two pulled-back
canonical divisors defined by each local frame on $Z$. A change
of frame adds the same pulled-back principal divisor on both
sides, so the differences glue to a global rational discrepancy
divisor $E_f$. Thus \eqref{eq:mixed-canonical-comparison} is a global
$\Q$-linear equivalence. Since $q$ is small, every prime of $X^+$
comes from a prime of $X$; hence $E_f$ is $a$-exceptional. On an
$a$-vertical curve $C$, the image $p(C)$ is a point or an
$f$-vertical curve, so
\[
 E_f\cdot C=p^*K_X\cdot C\le0.
\]
Negativity proves $E_f\ge0$.

If $E_f=0$, lift a nontrivial $f$-vertical curve to a curve on $W$.
Its degree against $p^*K_X$ is negative, whereas its degree against
$a^*K_{X^+}$ is nonnegative, a contradiction. If an $f$-exceptional
prime had zero coefficient, a general $a$-vertical curve through a
general point of its strict transform would not be contained in
$\Supp E_f$. Its $E_f$-degree would be nonnegative. Its image on
$X$ is nontrivial, since $p$ is an isomorphism at that general point,
so the same equality instead gives negative degree. This proves
strictness at every contracted original prime.

The comparison increases discrepancies. A valuation exceptional over $X$ has nonnegative discrepancy there.
A prime of $X$ contracted by the step has discrepancy zero on $X$
and acquires a positive coefficient in the comparison. These exhaust the exceptional
valuations over $X^+$, which is therefore canonical. Relative
Kawamata--Viehweg vanishing, in the analytic form
\cite[Theorem~5.2]{Fujino2022} with boundary and integral divisor
both zero, gives $R^if_*\OO_X=0$ for $i>0$.
A resolution of $X$ and Leray show that $Z$ has rational
singularities. If the original target is compact Kähler, apply the
construction to the entire original target, so that no restriction is needed.
A complex space projective over a compact Kähler space is compact
Kähler, giving the last assertion.
\end{proof}

\subsection{Supporting one negative analytic ray}

We use the generalized-pair terminology of \cite{HX2026,OAI034}.
The nef datum is represented by a nef class on a fixed smooth
Kähler carrier; its trace on a working model is the birational
pushforward of that class. All adjoint classes below are
Bott--Chern classes. The pushforward of a big carrier class is
modified big, and a modified-big Bott--Chern class on a Kähler
space is big \cite[Definition~2.8]{HX2026}.
The \emph{analytic cone} is the closed cone of classes generated
by positive closed currents of bidimension $(1,1)$.

\begin{lemma}\label{lem:canonical-detected}
Let $T$ be a canonical compact Kähler space. Let
$A=\{K_T\}+\beta_T$ be a gklt adjoint with globally nef carrier
data and big trace $\beta_T$. Suppose $R$ is an $A$-negative
extremal ray of the full analytic cone and $K_T\cdot R<0$.
Then $R$ has a projective contraction $f:T\to Z$ with normal
compact Kähler target and $-K_T$ relatively ample.
If $f$ is birational, it has the canonical replacement of
Lemma~\ref{lem:canonical-replacement} over all of $Z$.
\end{lemma}

\begin{proof}
We use the supporting-class argument of
\cite[proof of Lemma~3.6]{OAI034}.
The generalized cone theorem of Hacon--Li--Xie
\cite[Corollary~3.6]{HLX2026}, building on Hacon--P\u{a}un
\cite[Proposition~3.1 and Corollary~3.2]{HP2024}, expresses the
analytic cone as its $A$-nonnegative part and finitely many
$A$-negative rational rays; see also \cite[Theorem~1.3]{HX2026}.
We claim that $R$ is the null face of a nef supporting class
$\ell$. Normalize the analytic cone by a Kähler class. The cone generated by its $A$-nonnegative part and the
finitely many other negative rays is closed: its normalized slice
is the compact convex hull of the corresponding closed slices.
It does not contain the normalized point of $R$. Separate that
point from this cone and add the unique positive multiple of the
normalizing Kähler class which makes the value on $R$ zero. The
result is strictly positive on the complementary cone and zero on
$R$, as required by cone duality.

For $b\gg0$, the class $\omega_R=b\ell-A$ is Kähler. Indeed,
$-A$ is positive near $R$ on the compact normalized slice, and
$\ell$ has positive minimum on the complementary compact set.
Add the pullback of $\omega_R$ to the fixed nef carrier data.
The generalized discrepancies do not change; the new datum is nef
and its trace remains big. Theorem~1.4 of \cite{HX2026} applies
to the nef adjoint $A+\omega_R=b\ell$. It supplies a Moishezon
contraction to a normal compact Kähler space with
\[
 b\ell=f^*\gamma,\qquad\gamma\text{ Kähler on }Z.
\]
Its null analytic face is the single ray $R$.

Since $K_T$ is $\Q$-Cartier and $K_T\cdot R<0$,
\cite[Lemma~3.16]{OAI034} gives a relatively ample Cartier
multiple of $-K_T$, so $f$ is projective.
In this setting, positivity on the compact cone slice makes
$f^*\gamma-s\{K_T\}$ Kähler for small $s>0$; a local base
potential can be absorbed in a metric on a fixed Cartier multiple
of $-K_T$, proving relative ampleness of this multiple.
The birational assertion now follows from
Lemma~\ref{lem:canonical-replacement}, applied with $K=Z$.
\end{proof}

For scaling, we also need to transport Bott--Chern classes through
the step. The degree along the contracted ray determines the
exceptional correction; the remaining class descends to the base.

\begin{lemma}\label{lem:canonical-transport}
Consider a birational step obtained in
Lemma~\ref{lem:canonical-detected}. Every Bott--Chern class
$\theta$ on $T$ has a linear transform $\theta^+$ on $T^+$.
With the notation of \eqref{eq:mixed-canonical-comparison},
\begin{equation}\label{eq:canonical-linear-transport}
 p^*\theta-a^*\theta^+=c(\theta)\{E_f\},
 \qquad c(\theta)=\frac{\theta\cdot R}{K_T\cdot R}.
\end{equation}
The class $\theta^+$ represents the birational pushforward of
$\theta$. In particular, $\{K_T\}^+=\{K_{T^+}\}$.
\end{lemma}

\begin{proof}
The source and contraction base have rational singularities.
Relative vanishing gives $R^if_*\OO_T=0$ for $i>0$.
Every contracted curve has analytic class in $R$.
Consequently $\theta-c(\theta)\{K_T\}$ has degree zero on all
contracted curves. Degree-zero descent
\cite[Lemma~2.39]{HX2026} gives a unique class $\eta$ on $Z$ with
\[
 \theta-c(\theta)\{K_T\}=f^*\eta.
\]
Set $\theta^+=q^*\eta+c(\theta)\{K_{T^+}\}$ and use
\eqref{eq:mixed-canonical-comparison}. This proves linearity and
\eqref{eq:canonical-linear-transport}. Pushing its exceptional
difference to the new model proves the trace assertion.
\end{proof}

\section{Finite programs and divisorial negative parts}
\label{sec:canonical-program}

We first construct canonical programs with scaling. A termination
criterion, expressed in terms of divisorial negative parts, then
produces the finite program used in
Section~\ref{sec:kollar-family}. The same criterion will also apply
to the logarithmic program of Section~\ref{sec:finite}.

\subsection{Divisorial negative parts}
\label{subsec:negative-parts}

For a pseudoeffective Bott--Chern class $\alpha$ on a smooth compact
Kähler manifold, let
\[
 \nu_D(\alpha)=\lim_{\varepsilon\downarrow0}
  \inf\{\nu_D(T):T\in\alpha,\ T\ge-\varepsilon\omega\},
 \qquad N(\alpha)=\sum_D\nu_D(\alpha)D,
\]
where $\omega$ is a fixed Kähler form. This definition is independent
of $\omega$, and $N(\alpha)$ is an effective divisor with finite support
\cite[Sections~2--3 and Theorem~3.12]{Boucksom}. For a holomorphic
$\Q$-line bundle $J$, write $N(J)=N(\{J\})$. For a holomorphic $\Q$-line bundle on a normal compact Kähler
space, a divisorial multiplicity is computed after pullback to a
smooth Kähler resolution carrying that divisor. On a given model,
$N(J)$ denotes the sum over the prime divisors of that model. For effective divisors $N,B$, the notation
$N\wedge B=0$ means that they have no common prime component.

We use homogeneity, monotonicity under addition of a nef class, and
exceptional translation: if $q:W\to Y$ is a modification, $\alpha$
is pseudoeffective, and $E\ge0$ is $q$-exceptional, then
\begin{equation}\label{eq:exceptional-translation}
 N(q^*\alpha+\{E\})=N(q^*\alpha)+E.
\end{equation}
The normal-space formulation and its compatibility with higher
resolutions are recorded in \cite[Lemmas~2.2--2.4]{OAI034}.

\subsection{Canonical scaling}

A \emph{positive truncation} of a scaling program consists of the
steps whose scaling thresholds are at least a fixed number
$\varepsilon>0$. The scaling parameter starts at $1$; the initial
Kähler class is chosen so that the adjoint at $1$ is Kähler.

\begin{lemma}\label{lem:canonical-scaling}
Let $X$ be a canonical compact Kähler space. Choose a sufficiently
large Kähler class $H$ with $\{K_X\}+H$ Kähler. There is a $K_X$-negative program with scaling of $H$, possibly
infinite, whose birational maps
are the canonical steps of Section~\ref{sec:canonical-steps}, and
all working spaces are canonical and compact Kähler. It stops
when the canonical class is nef or upon reaching a projective
$K$-negative fiber-type contraction.

Write $\alpha_i(s)=\{K_{X_i}\}+sH_i$ for the linear traces.
Every positive truncation of the program is finite. If the program
is infinite, its thresholds tend to zero. At a wall $s_i>0$,
the scaled pullbacks agree. On a common resolution
$p_i:W_i\to X_i$, $a_i:W_i\to X_{i+1}$,
\begin{equation}\label{eq:nonfactorial-wall}
 p_i^*\alpha_i(s)-a_i^*\alpha_{i+1}(s)
 =\left(1-\frac{s}{s_i}\right)\{E_i\},
 \qquad E_i\ge0,\quad E_i\ne0.
\end{equation}
If $K_X$ is not pseudoeffective, this program is finite and ends
with a projective $K$-negative fiber-type contraction.
\end{lemma}

\begin{proof}
Start with $X_0=X$, $H_0=H$, and upper parameter $u=1$.
Suppose inductively that a finite prefix has reached $X_i$, that
$\alpha_i(u)$ is nef, and that for $0\le s\le u$ its comparison
with the original scaled class is effective and exceptional over
$X_i$. On a common smooth carrier, the nef datum remains the
pullback of $sH$ from $X$. Its trace on $X_i$ is exactly $sH_i$
by Lemma~\ref{lem:canonical-transport}. Since the pullback of $H$
is nef and big, this trace is modified big for $s>0$. It is also
a Bott--Chern class on the Kähler space $X_i$, so it is big
\cite[Definition~2.8]{HX2026}. On the common carrier, the effective comparison subtracts an
effective divisor from the original generalized boundary. Thus
$(X_i,0+s\mathbf H)$ is gklt for $0\le s\le u$. Its ordinary
boundary is zero, since all preceding steps are nonextracting.

If $K_{X_i}$ is nef, stop. Otherwise let
\[
 \tau=\min\{s\in[0,u]:\alpha_i(s)\text{ is nef}\}>0,
 \qquad 0<t_0<\tau.
\]
Apply the generalized cone theorem
\cite[Corollary~3.6]{HLX2026} to the gklt big-trace adjoint
$A=\alpha_i(t_0)$. Its negative cone has finitely many extremal
rays. On the $A$-nonnegative part, all $\alpha_i(s)$ with
$t_0\le s\le u$ are nonnegative, being convex combinations of
$A$ and $\alpha_i(u)$. Thus $\tau$ is determined by the finitely
many negative-ray inequalities. One such ray $R$ satisfies
\[
 \alpha_i(\tau)\cdot R=0,
 \qquad\alpha_i(t_0)\cdot R<0.
\]
It follows that $H_i\cdot R>0$ and
$K_{X_i}\cdot R=-\tau H_i\cdot R<0$.
Lemma~\ref{lem:canonical-detected} constructs the contraction of
this single ray, using a new Kähler supporting direction on $X_i$.
A fiber-type contraction ends the program. In the birational case,
write its canonical replacement as
$X_i\xrightarrow{f_i}Z_i\xleftarrow{q_i}X_{i+1}$.

The wall annihilates $R$, so degree-zero descent and
Lemma~\ref{lem:canonical-transport} give
\[
 \alpha_i(\tau)=f_i^*\delta,
 \qquad\alpha_{i+1}(\tau)=q_i^*\delta.
\]
The class $\delta$ is nef: nefness descends under a surjective
map of normal compact Kähler spaces, by the general pullback
criterion recalled in \cite[Section~2.4]{HX2026}. Therefore the
new wall is nef. The unscaled canonical comparison now gives
\eqref{eq:nonfactorial-wall}, with $s_i=\tau$.
For $s\le\tau$ it preserves the effective-comparison and gklt
assertions used above, so the construction can be repeated.
The finiteness argument below also bounds the number of steps at
any fixed positive threshold.

Fix $\varepsilon>0$ and a resolution $\nu:W\to X$ of the
original model. The compact polytope
\[
 \{(0,s\nu^*H):\varepsilon\le s\le1\}
\]
has globally nef carrier data, gklt original pairs, and big traces.
The finite marked weak-model theorem
\cite[Theorem~5.15(2)]{HX2026} applies in the category of
Definition~2.15 of that paper. The nef wall and the composed effective comparison make every
working model with threshold at least $\varepsilon$ a marked weak
model in that theorem. A marked model cannot recur:
every nontrivial step has a nonzero effective unscaled canonical
comparison, these comparisons add after pullback, and recurrence
would give identical canonical discrepancies. There are therefore only finitely many steps with threshold at
least $\varepsilon$. In particular, an infinite program has
thresholds tending to zero.

If $K_X$ is not pseudoeffective, let $\epsilon>0$ be its
pseudoeffective threshold in the direction $H$. Every nef working wall, together with its effective
exceptional comparison, makes the original scaled class
pseudoeffective. Thus every threshold is at least $\epsilon$.
Positive-truncation finiteness makes the program finite. A nef
endpoint would make the unscaled original canonical class
pseudoeffective by its effective comparison, a contradiction.
The final contraction is therefore of fiber type.
\end{proof}

\subsection{A termination criterion}

The next criterion abstracts the affine comparison argument of
\cite[proof of Theorem~3.13]{OAI034}. First, a finite prefix
removes all divisorial negative parts on the working model. A new
scaling direction then gives a finite continuation to a nef model.

\begin{proposition}[Termination from one nef interval]
\label{prop:termination-criterion}
Let $\mathscr C$ be a class of normal compact Kähler models $X$
equipped with holomorphic $\Q$-line bundles $J$, closed under the chosen
negative steps. Suppose the following properties hold for every
pseudoeffective member $(X,J)$ of $\mathscr C$.
\begin{enumerate}[label=\textup{(\roman*)},leftmargin=*]
\item For every sufficiently large Kähler class $H$, a $J$-negative
program with scaling of $H$ either reaches a nef model or continues
indefinitely. Its maps are nonextracting, its transforms $J_i$
are pseudoeffective, and every positive truncation is finite.
The thresholds are positive and nonincreasing. Each scaled class
$\{J_i\}+sH_i$ is nef between its incoming and outgoing thresholds;
the classes $H_i$ transform linearly, and the scaled pullbacks
agree at each wall $s_i$.
On a common resolution of each step,
\[
 p_i^*J_i\sim_{\Q}q_i^*J_{i+1}+E_i,
 \qquad E_i\ge0,\quad q_{i*}E_i=0,\quad E_i\ne0,
\]
with positive coefficient at every prime contracted from $X_i$.
\item If $N(J_T)=0$ on a working model $T$, then for each
sufficiently large Kähler class $A$ there are a compact Kähler
space $V$ small bimeromorphic with $T$, transformed classes
$\{J_V\},A_V$, and $\delta>0$ such that
$\{J_V\}+sA_V$ is nef for $0\le s\le\delta$.
The two pullback differences on a common smooth resolution are
represented by real divisors exceptional over both models.
\end{enumerate}
Then every pseudoeffective member $(X,J)$ admits a finite
$J$-negative program starting on $X$ and ending at a nef model. If, on a resolution,
each $J_i$ has a decomposition
$r^*J_i\sim_{\Q}P_i+N(r^*J_i)$ with $P_i$ semiample, the endpoint
is semiample.
\end{proposition}

\begin{proof}
\emph{Scaled comparisons.}
Wall agreement and linearity give
\begin{equation}\label{eq:scaling-comparison}
 p_i^*(\{J_i\}+sH_i)-q_i^*(\{J_{i+1}\}+sH_{i+1})
 =\left(1-\frac{s}{s_i}\right)\{E_i\}.
\end{equation}
On a nef working interval, compose these formulas. All previous
thresholds are at least the parameter, so the resulting comparison
is effective and exceptional over the working model. By
\eqref{eq:exceptional-translation}, it is the full negative part
of the original scaled class on the common resolution.

\smallskip\noindent\emph{Removing divisorial negative parts.}
For the unscaled comparison, pseudoeffectivity and
\eqref{eq:exceptional-translation} place each contracted original
prime in $\Supp N(J)$. This support is finite and
nonextraction prevents a prime from reappearing. Thus only finitely
many steps contract divisors.

Suppose the program is infinite. Its positive thresholds tend to
zero, by positive-truncation finiteness. Choose a model $T$ after
the last divisorial step and after a finite equal-threshold block,
with a nonempty nef working interval. For a parameter $\lambda_0$
in this interval, put $Q=\{J_T\}+\lambda_0H_T$, which is nef.
All later maps are small. For parameters $\lambda$ in the later nef working intervals, the
exceptional comparison has zero coefficient at each prime of $T$.
Hence the minimal multiplicities of $\{J_T\}+\lambda H_T$ at
those primes vanish. Such parameters tend to zero.

To pass to the limit, write
\begin{equation}\label{eq:nef-convex-limit}
 \{J_T\}+\lambda H_T
 =\left(1-\frac{\lambda}{\lambda_0}\right)\{J_T\}
   +\frac{\lambda}{\lambda_0}Q.
\end{equation}
On a fixed smooth Kähler resolution $r:W\to T$, choose a Kähler
class $\Omega$ and $c>0$ with $c\Omega-r^*Q$ nef. Then
\[
 \nu_D(r^*\{J_T\}+c\eta\Omega)
 \le\nu_D(r^*\{J_T\}+\eta r^*Q)
 \le\nu_D(r^*\{J_T\})
\]
for every prime $D$ on $W$. By the definition of minimal multiplicity, the outer terms have
the same limit as $\eta\downarrow0$. Homogeneity and
\eqref{eq:nef-convex-limit} therefore give convergence to the
multiplicities of $J_T$. It follows that $N(J_T)=0$ on $T$.

\smallskip\noindent\emph{A finite continuation.}
Restart scaling on $T$ with a sufficiently large Kähler class $A$.
The preceding strict comparison excludes divisorial steps. Choose
$V,\delta$ as in \textup{(ii)}. On a common smooth resolution
$p:W\to T$, $q:W\to V$, write
\begin{equation}\label{eq:affine-negative-part}
 p^*(\{J_T\}+sA)
 =q^*(\{J_V\}+sA_V)+\{F(s)\},
 \qquad 0\le s\le\delta.
\end{equation}
The divisor $F(s)$ is affine and exceptional over both models.
Its negative is $p$-nef, so exceptional negativity gives $F(s)\ge0$.
Equation~\eqref{eq:exceptional-translation} now yields
\begin{equation}\label{eq:single-affine-negative}
 N\bigl(p^*(\{J_T\}+sA)\bigr)=F(s)
 \qquad(0\le s\le\delta).
\end{equation}
This equality persists on higher resolutions.

If the restarted program were infinite, group the finitely many
steps at each positive threshold into a block. Its effective
unscaled comparisons add to a nonzero effective divisor, and
\eqref{eq:scaling-comparison} still holds for the block. The resulting distinct walls would tend to zero. Choose a
wall $s_i\in(0,\delta)$ whose two adjacent working intervals also
lie in $(0,\delta)$, and resolve the corresponding finite prefix
with $T$ and $V$. The nef comparisons on the two intervals compute
the negative part of the same original class. At a prime having
coefficient $e_i>0$ in the block comparison, the two affine slopes
differ by $e_i/s_i$. They cannot both agree with the single affine
coefficient in \eqref{eq:single-affine-negative}. Thus the
restarted program terminates. Joining it to the finite initial
prefix gives the asserted program from $X$.

\smallskip\noindent\emph{Semiampleness.}
At a nef endpoint, the good divisorial decomposition has zero
negative part on a resolution. The pulled-back holomorphic
$\Q$-line bundle is consequently semiample. Clear its Cartier index; normality
and the projection formula identify its sections with those on the
endpoint. Generation upstairs then implies generation downstairs.
\end{proof}

\subsection{The finite canonical program}

\begin{proposition}\label{prop:canonical-finite}
Let $X$ be a normal compact Kähler space with canonical singularities.
There is a finite program of canonical negative steps starting on
$X$, with canonical compact Kähler working spaces. If $K_X$ is
pseudoeffective, its endpoint has semiample canonical divisor.
Otherwise it ends with a projective $K$-negative fiber-type contraction.
\end{proposition}

\begin{proof}
The non-pseudoeffective case is Lemma~\ref{lem:canonical-scaling}.
In the pseudoeffective case, Proposition~\ref{prop:decomposition-input}
gives on a smooth compact Kähler modification $r:W\to X$ a
$\Q$-linear equivalence of holomorphic $\Q$-line bundles
\begin{equation}\label{eq:canonical-good-decomposition}
 r^*K_X\sim_{\Q}P+N(r^*K_X),\qquad P\text{ semiample}.
\end{equation}
Thus some pluricanonical power has a nonzero section. Canonical
bimeromorphic invariance (Lemma~\ref{lem:birational}) preserves
this property on every working model. It excludes a negative
fiber-type contraction: after raising the section to a
Cartier-index-clearing power, its nonzero restriction to a general
positive-dimensional projective fiber would lie in an antiample line.
Every contraction target is compact, so
Lemma~\ref{lem:canonical-replacement} applies over the entire target.
Together with Lemma~\ref{lem:canonical-scaling}, this verifies condition~\textup{(i)} of
Proposition~\ref{prop:termination-criterion} for the class of
canonical working models.

For condition~\textup{(ii)}, apply
Proposition~\ref{prop:interval-input} to the klt pair $(T,0)$.
Proposition~\ref{prop:termination-criterion} now gives a finite
program and, by \eqref{eq:canonical-good-decomposition}, a
semiample endpoint.
\end{proof}

\section{Canonical invariance in analytic families}
\label{sec:kollar-family}

We follow Kollár's construction \cite[paragraph~8]{Kollar2021}:
extend the negative contraction and recover the positive side from
its relative canonical algebra. Fujino's analytic finite-generation
theorem \cite[Theorem~1.18]{Fujino2022} supplies that algebra.
Pluricanonical residue identifies its central quotient and hence
the central positive model.

\Needspace{19\baselineskip}
\subsection{Deformation of canonical steps}

\begin{lemma}[Deformation of a canonical MMP step]
\label{lem:deform-mixed}
Let $g:\mathcal X\to\Delta$ be proper and flat with normal connected
canonical central fiber. Suppose that
\[
 X_0\xrightarrow{f_0}Z_0\xleftarrow{q_0}Y_0
\]
is a diagram of projective bimeromorphic contractions of normal
compact complex spaces, with $-K_{X_0}$ relatively ample for $f_0$,
$q_0$ small, and $K_{Y_0}$ relatively ample for $q_0$.
Assume $Y_0$ is canonical. After shrinking the disk there is a
relative diagram
\[
 \mathcal X\xrightarrow{f}\mathcal Z\xleftarrow{q}\mathcal Y
\]
having this central diagram, with $-K_{\mathcal X/\Delta}$
$f$-ample and $K_{\mathcal Y/\Delta}$ $q$-ample. The families are
proper and flat, and the nearby fibers of $\mathcal X$ and
$\mathcal Y$, as well as their total spaces, are canonical.
The induced correspondences $X_t\dashrightarrow Y_t$ are
bimeromorphic on every nearby fiber.
\end{lemma}

\begin{proof}
\emph{Extending the contraction.}
By \cite[Lemma~10 and Remark~10.3]{Kollar2021}, canonical
singularities of the compact central fiber give, after shrinking,
normal canonical total space, a $\Q$-Cartier relative canonical
divisor, and canonical nearby fibers. Relative Kawamata--Viehweg
vanishing \cite[Theorem~5.2]{Fujino2022}, applied to the
projective $K$-negative contraction with $D=0$, gives
$R^1(f_0)_*\OO_{X_0}=0$. The contraction-extension theorem
\cite[Proposition~11.4]{KM1992}, in the form stated in
\cite[Theorem~4.7]{HLR2026}, extends $f_0$ to a contraction
$f:\mathcal X\to\mathcal Z$ over a smaller disk, with proper
flat target and central fiber $Z_0$.
The identity $f_*\OO_{\mathcal X}=\OO_{\mathcal Z}$ gives a normal
target. Its parameter is a nonzerodivisor, so it is flat over the
disk; the parameter sequence and the central $R^1$ vanishing
identify its central structure sheaf with $\OO_{Z_0}$.

Take a relative Cartier multiple of $-K_{\mathcal X/\Delta}$.
Its central restriction is $f_0$-ample. Openness of relative
ampleness \cite[Proposition~3.5]{HLR2026} and compactness make
this line $f$-ample after shrinking. Thus $f$ is projective. It is bimeromorphic because its source
and target have the same dimension and
$f_*\OO_{\mathcal X}=\OO_{\mathcal Z}$.

\smallskip\noindent\emph{The positive model and its central fiber.}
Apply Lemma~\ref{lem:canonical-replacement} to $f$ with the compact
subset $Z_0\subset\mathcal Z$. It gives the canonical replacement over
an open neighborhood $U$ of $Z_0$. Write $\pi:\mathcal Z\to\Delta$
for the proper target family. The set $\pi(\mathcal Z\setminus U)$ is
closed and does not contain $0$. Choose a smaller disk
$\Delta_\varepsilon$ disjoint from this set, so that
$\pi^{-1}(\Delta_\varepsilon)\subset U$. Restrict the diagram to
this disk and retain the notation $\Delta$ for it. Denote the resulting
canonical replacement by $q:\mathcal Y\to\mathcal Z$.
The relative canonical algebra is locally finitely generated;
$\mathcal Y$ is canonical, $q$ is small, and
$K_{\mathcal Y/\Delta}$ is $q$-ample.
For the residue calculation, put
\begin{equation}\label{eq:family-canonical-algebra}
 \mathscr R=\bigoplus_{m\ge0}
 \omega_{\mathcal Z}^{[m]}(mZ_0).
\end{equation}
Twisting by the Cartier divisor $Z_0$ does not change the Proj.

Choose a projective
resolution $\tau:\mathcal W\to\mathcal X$ such that the strict
transform $W_0$ of $X_0$ is smooth. Put $h=f\tau$.
The smooth divisor $W_0$ is nonexceptional for the projective
bimeromorphic map $h$. Nakayama's theorem
\cite[VI.3.8]{Nakayama2004}, in the precise analytic form of
\cite[Theorem~9]{Kollar2021}, gives
\begin{equation}\label{eq:nakayama-relative-surjection}
 h_*\omega_{\mathcal W}^{\otimes m}(mW_0)
 \twoheadrightarrow
 (h|_{W_0})_*\omega_{W_0}^{\otimes m}.
\end{equation}
Canonicality of $X_0$ and the negative-side algebra identity identify
the right side with $\omega_{Z_0}^{[m]}$. Codimension-one regularity identifies the left side with a
subsheaf of $\omega_{\mathcal Z}^{[m]}(mZ_0)$. The flat morphism $\mathcal Z\to\Delta$ is smooth along
$(Z_0)_{\reg}$. Pluricanonical residue is therefore defined there and extends to
$\omega_{Z_0}^{[m]}$ by reflexivity. It agrees with the residue in
\eqref{eq:nakayama-relative-surjection}. We obtain a surjection
\[
 \omega_{\mathcal Z}^{[m]}(mZ_0)
 \twoheadrightarrow\omega_{Z_0}^{[m]}.
\]
Its kernel is exactly $t\omega_{\mathcal Z}^{[m]}(mZ_0)$.
Indeed, a section of zero residue is divisible by $t$ at the generic
point of $Z_0$. At every other codimension-one point, $t$ is a unit.
The quotient by $t$ consequently remains in the same reflexive
sheaf by normality. Conversely, a multiple of $t$ has zero residue.
The restrictions are multiplicative, hence
\begin{equation}\label{eq:canonical-algebra-base-change}
 \mathscr R/t\mathscr R
 \simeq\bigoplus_{m\ge0}\omega_{Z_0}^{[m]}.
\end{equation}

Relative Proj commutes with the quotient by $t$. Smallness of
$q_0$ and $q_0$-ampleness of $K_{Y_0}$ identify the Proj of the
right-hand algebra in \eqref{eq:canonical-algebra-base-change}
with $Y_0$. Thus the actual central fiber of $\mathcal Y$ is
the prescribed $Y_0$. The parameter is a nonzerodivisor on the
normal bimeromorphic total space, so $\mathcal Y\to\Delta$ is
flat; it is proper by construction. Canonicality of its nearby
fibers follows from \cite[Lemma~10]{Kollar2021} after shrinking.

\smallskip\noindent\emph{Fiberwise bimeromorphicity.}
Both central arrows are isomorphisms over a common nonempty open
subset of $Z_0$. Lemma~\ref{lem:open-isomorphism}
extends this common isomorphism locus to a total open set.
Openness of $\mathcal X\to\Delta$ makes this set meet every
sufficiently nearby fiber. Upper semicontinuity of $h^0(\OO)$ makes the nearby normal
fibers $X_t$ and $Y_t$ connected, hence irreducible. Their identified nonempty open subsets therefore
give the asserted bimeromorphic correspondences.
\end{proof}

\begin{lemma}[Deformation of the final negative fiber contraction]
\label{lem:deform-mori}
Let $\mathcal M\to\Delta$ be proper and flat, with connected canonical
central fiber, and let $h_0:M_0\to Z_0$ be a projective contraction
such that $\dim Z_0<\dim M_0$ and $-K_{M_0}$ is $h_0$-ample.
After shrinking, all reflexive plurigenera of all fibers of
$\mathcal M$ are zero, on one disk independent of the degree.
\end{lemma}

\begin{proof}
Canonical openness gives a canonical total space and canonical
nearby fibers after shrinking. Upper semicontinuity of
$h^0(\OO_{M_t})$ makes these fibers connected.
Relative vanishing \cite[Theorem~5.2]{Fujino2022} gives
$R^1(h_0)_*\OO_{M_0}=0$, and
\cite[Proposition~11.4]{KM1992} extends the contraction.
This is the zero-boundary case of
\cite[Proposition~4.9]{HLR2026}. The target is proper over the disk,
and its fiber dimension is at most $\dim Z_0$ near zero.
One common Cartier index $q$ can be chosen for
$K_{\mathcal M/\Delta}$. Openness of relative ampleness makes
\[
 A=\OO_{\mathcal M}(-qK_{\mathcal M/\Delta})
\]
relatively ample for the extended contraction. Since
$\dim Z_t\le\dim Z_0<\dim M_t$, the contraction has
positive-dimensional fibers on a dense open subset of each nearby
$M_t$.

Suppose $s\in H^0(M_t,\omega_{M_t}^{[m]})$ is nonzero. Its
$q$-th reflexive power is a nonzero section of $A_t^{-m}$.
Choose a positive-dimensional projective fiber meeting its
nonvanishing locus, and then an integral curve in that fiber
through a nonvanishing point. The restricted section is nonzero,
whereas $A_t^{-m}$ has strictly negative degree on that curve.
This is impossible. Since $A$ and the neighborhood were fixed before choosing $m$,
the vanishing holds in every degree on that neighborhood.
\end{proof}

\subsection{Families over a disk}

\begin{proposition}[The disk case]\label{prop:canonical-disk}
The assertions of Theorem~\ref{thm:canonical-central} hold when
the base is a disk.
\end{proposition}

\begin{proof}
\emph{Reduction to the endpoint.}
Shrink using \cite[Lemma~10 and Remark~10.3]{Kollar2021}.
The total space and all fibers are normal canonical, and the
relative canonical divisor is $\Q$-Cartier. Upper semicontinuity of
$h^0(\OO_{X_t})$ makes the nearby fibers connected, hence irreducible.
The reflexive sheaves
\[
 \mathcal F_m=\OO_{\mathcal X}(mK_{\mathcal X/\Delta})
\]
are flat over the disk and their actual fiber restrictions are
$\omega_{X_t}^{[m]}$ by Lemma~\ref{lem:restrictions}.

Proposition~\ref{prop:canonical-finite} gives a finite
birational program
\[
 X_0=X_0^0\dashrightarrow X_0^1\dashrightarrow\cdots
 \dashrightarrow X_0^N=M_0
\]
whose endpoint either has semiample canonical divisor or carries
a projective $K$-negative fiber-type contraction.
All central working spaces are canonical and compact Kähler.
Lemma~\ref{lem:deform-mixed} extends the birational list over
one smaller disk. Applying Lemma~\ref{lem:birational} with zero
boundary to every fiber gives
\begin{equation}\label{eq:canonical-birational-invariance}
 P_m(X_t)=P_m(M_t)\qquad(t\in\Delta,\ m\ge1).
\end{equation}

In the semiample case, $M_0$ has a compact Kähler resolution,
so Corollary~\ref{cor:uniform-endpoint} makes every $P_m(M_t)$
constant on one disk. In the fiber-type case,
Lemma~\ref{lem:deform-mori} makes all these plurigenera zero on
one disk. Thus in both cases, simultaneously for every $m\ge1$,
\[
 P_m(X_t)=P_m(M_t)=P_m(M_0)=P_m(X_0).
\]

\smallskip\noindent\emph{Direct images and section extension.}
Apply Grauert's theorem to the flat coherent sheaf
$\mathcal F_m$ for each $m$ on the already fixed disk.
Each $f_*\mathcal F_m$ is locally free and has holomorphic base
change. Surjectivity on the central fiber gives extension of every
central pluricanonical section. Since the disk is Stein, Cartan's theorem~A extends every such
vector to a global section of its direct-image bundle.
\end{proof}

\begin{remark}[Projective central fibers]
For a projective central fiber, the finite program can also be obtained
from the nonfactorial termination theorem of Hashizume--Hu
\cite[Theorem~1.7]{HH2020}. If $K_{X_0}$ is
pseudoeffective, \cite[Proposition~2.10]{OAI034} gives a good
minimal model. The theorem of Hashizume--Hu supplies a terminating
negative MMP beginning on the given $X_0$, without a factoriality
assumption. Its endpoint is
good by \cite[Lemma~2.10]{TX2024}. The mixed steps described in
\cite[Remark~2.15]{TX2024} are precisely the diagrams of
\cite[Definition~7]{Kollar2021}. Hence
\cite[Theorem~2]{Kollar2021} extends the entire finite program.
The uniform endpoint theorem and canonical bimeromorphic invariance
give the disk assertion of Proposition~\ref{prop:canonical-disk}.
If $K_{X_0}$ is not pseudoeffective,
\cite[Theorem~1.7]{HH2020} instead supplies a Mori program,
and Lemma~\ref{lem:deform-mori} handles all degrees together.
\end{remark}

\subsection{General analytic bases}

We use the following disk criterion for the full sequence of
plurigenera; compare \cite[paragraph~13]{Kollar2021}.

\begin{lemma}[Detection by disks]\label{lem:disk-detection}
Let $S$ be a complex analytic space, possibly nonreduced, let
$\Lambda$ be a set, and let $u:|S|\to\Lambda$ be a function.
If $u\circ\gamma$ is constant for every holomorphic map
$\gamma:\Delta\to S$, then $u$ is locally constant.
\end{lemma}

\begin{proof}
The assertion is local on $S$. Reduction preserves the underlying
topological space, and every map from a disk factors through
$S_{\mathrm{red}}$. Take a proper surjective resolution
$\pi:T\to S_{\mathrm{red}}$
\cite[Theorem~13.3]{BM1997}, using the disjoint union of resolutions
of the local irreducible components when necessary.
Any two points in a coordinate polydisk of $T$ can be joined by a
finite chain of coordinate disks. The hypothesis therefore makes
$u\circ\pi$ constant on every such polydisk, hence locally constant.
For $s_0\in S$, the set
\[
 W=\{t\in T:u(\pi(t))=u(s_0)\}
\]
is open and contains $\pi^{-1}(s_0)$. Since $\pi$ is proper,
\[
 U=S\setminus\pi(T\setminus W)
\]
is an open neighborhood of $s_0$. Surjectivity gives
$u(s)=u(s_0)$ for every $s\in U$.
\end{proof}

\begin{proof}[Proof of Theorem~\ref{thm:canonical-central}]
Proposition~\ref{prop:canonical-disk} gives the assertions over a disk.
For the numerical assertion over an arbitrary base, replace $S$
by $S_{\mathrm{red}}$ and pull back the family. Properness,
flatness and all point fibers are preserved.

Nakayama's stability theorem \cite[VI.5.2]{Nakayama2004}, in the
analytic form of \cite[Lemma~10]{Kollar2021}, and properness give
a neighborhood of $0$ on which every fiber is canonical.
Canonical singularities are rational. By Bingener's deformation
theorem \cite[Theorem~6.3]{Bingener1983}, in the form of
\cite[Corollary~3.3.6]{Kirschner2012}, all these fibers are
Kähler after a further restriction. Openness of flat maps and upper
semicontinuity of $h^0(X_s,\OO_{X_s})$ also make the fibers nonempty
and connected. Denote the resulting neighborhood by $V$ and set
\[
 \Phi:V\longrightarrow\Z_{\ge0}^{\Z_{>0}},\qquad
 \Phi(s)=\bigl(P_m(X_s)\bigr)_{m\ge1}.
\]

For a holomorphic map $\gamma:\Delta\to V$, the pullback family
is proper and flat, and every fiber satisfies the central-fiber
hypotheses of Proposition~\ref{prop:canonical-disk}. Applying that
proposition at each point of the disk shows that $\Phi\circ\gamma$
is locally constant, hence constant because the disk is connected.
Lemma~\ref{lem:disk-detection} implies
that $\Phi$ is locally constant, giving one neighborhood of $0$ for all
degrees. Reduction preserves both point fibers and underlying open
sets, so this is also the required neighborhood on the original base.
\end{proof}

\section{Canonical algebras and deformation consequences}
\label{sec:consequences}

\subsection{Canonical and adjoint algebras}

The semiample family constructed in the proof of
Proposition~\ref{prop:canonical-disk} determines the relative canonical
algebra. Birational comparison identifies its sections and their
products with those of the original family.

\begin{proof}[Proof of Theorem~\ref{thm:canonical-algebra}]
Under hypothesis \textup{(a)}, suppose first that
$\kappa(X_0)=-\infty$. Theorem~\ref{thm:canonical-central}
gives $P_m(X_t)=0$ for every $m>0$ on one disk $U$. The fibers are
connected and normal, so $f_*\OO_{\mathcal X}=\OO_U$.
Thus $\mathcal R(f)=\OO_U$ is concentrated in degree zero and
$\Proj_U\mathcal R(f)$ is empty.

In the remaining cases $\kappa(X_0)\ge0$. Under hypothesis \textup{(a)},
retain the proper flat family $g:\mathcal Y\to U$ obtained at the
semiample endpoint in the proof of
Proposition~\ref{prop:canonical-disk}. Its total space and fibers are
canonical, and the correspondence with $\mathcal X$ is bimeromorphic
on every fiber. Under hypothesis \textup{(b)}, take
$\mathcal Y=\mathcal X$ and apply
Corollary~\ref{cor:uniform-endpoint}.

Choose a Cartier multiple $kK_{\mathcal Y/U}$ whose central
restriction is globally generated. Lemma~\ref{lem:uniform} extends
a generating set and makes this multiple relatively generated on
one disk. Hence $K_{\mathcal Y/U}$ is relatively semiample, proving
\textup{(iii)}.
The sheaves $\OO_{\mathcal Y}(mK_{\mathcal Y/U})$ form a graded
algebra under multiplication of meromorphic pluricanonical forms.
Lemma~\ref{lem:restrictions} identifies their fiber restrictions and
gives their flatness. Tensoring by $\OO_{\mathcal Y}(kK_{\mathcal Y/U})$
gives the residue-class identities required by
Proposition~\ref{prop:endpoint-graded-algebra}. That proposition proves
\textup{(i)} and \textup{(ii)} for $g$.

Under \textup{(a)}, Lemma~\ref{lem:birational} on the total spaces
identifies the direct-image canonical algebras of $\mathcal X$ and
$\mathcal Y$. The identifications preserve multiplication because
they identify meromorphic pluricanonical forms. Applying the same
lemma on each fiber gives the stated fiber algebras.
The disk assertion of Theorem~\ref{thm:canonical-central} gives base change for the original
sheaves $\mathcal F_m$, and naturality of multiplication gives the
graded base-change isomorphism. This proves \textup{(i)} and
\textup{(ii)} for $f$.
\end{proof}

For $\kappa(X_0)=\dim X_0$, the relative Proj is the relative
canonical model. For $0\le\kappa(X_0)<\dim X_0$, it is the relative
Iitaka model, with fibers of dimension $\kappa(X_0)$.

\begin{corollary}[The admissible adjoint algebra]
\label{cor:relative-adjoint-algebra}
Under the assumptions of Theorem~\ref{thm:endpoint}, let $d>0$ be the
least integer for which $d\mathcal B$ and $dB$ are integral.
After restriction to a disk $U$ about $0$, $\mathcal A$ is relatively
semiample and
\[
 \mathcal R_d(g,\mathcal A)
 :=\bigoplus_{j\ge0}g_*\OO_{\mathcal Y}(jd\mathcal A)
\]
is a finitely generated flat $\OO_U$-algebra with locally free graded
pieces. The natural graded base-change map
\eqref{eq:graded-holomorphic-base-change} is an isomorphism for every
holomorphic base change. Its relative analytic Proj is normal, projective and
flat over $U$, and its formation also commutes with base change.
\end{corollary}

\begin{proof}
Corollary~\ref{cor:uniform-log-endpoint} gives relative semiampleness
and a common disk for all admissible degrees. The sheaves
$\mathcal F_j=\OO_{\mathcal Y}(jd\mathcal A)$ form a graded algebra,
and their residue-class identities respect multiplication.
Proposition~\ref{prop:endpoint-graded-algebra} applies and proves
the remaining assertions.
\end{proof}

\subsection{Higher pluricanonical cohomology}

Nakayama's theorem gives local freeness of pluricanonical higher
direct images for a proper family over a disk with klt total space,
relatively semiample canonical divisor and class-$\mathcal C$ fiber
components \cite[Corollary~3.15]{Nakayama1987};
see \cite[Proposition~4.6]{LiRaoWang2026}. For smooth families, the
following result assumes semiampleness and class $\mathcal C$ on
the central fiber.

The fixed-degree argument below combines Levine's cyclic-cover
construction with the eigenspace argument credited there to
Shepherd-Barron \cite[Corollary~1.8]{Levine1983}.
Lemma~\ref{lem:uniform-higher} then makes the neighborhood uniform
in both the canonical and cohomological degrees.

\begin{corollary}[Higher pluricanonical cohomology]
\label{cor:smooth-semiample-higher}
Let $f:\mathcal X\to\Delta$ be a proper holomorphic submersion of complex manifolds
with connected fibers. Suppose that $X_0$ belongs to Fujiki's
class $\mathcal C$ and that $K_{X_0}$ is semiample.
There is a disk $U$ about $0$ such that
\[
 h^i(X_t,K_{X_t}^{\otimes m})
 =h^i(X_0,K_{X_0}^{\otimes m})
 \qquad(t\in U,\ i\ge0,\ m\in\Z).
\]
For every $i\ge0$ and $m\in\Z$, the sheaf
$R^if_*K_{\mathcal X/U}^{\otimes m}$ is locally free on $U$ and
commutes with holomorphic base change.
\end{corollary}

\begin{proof}
Corollary~\ref{cor:uniform-endpoint} gives a relatively generated
power of $K_{\mathcal X/\Delta}$ on a disk about $0$.
Set $n=\dim X_0$; cohomology in degrees $i>n$ vanishes.
We first prove constancy for a fixed $m\ge1$ and all $0\le i\le n$.
Choose a sufficiently divisible $q\ge m$
such that $K_{X_0}^{\otimes q}$ is globally generated. A general
section has a smooth, possibly empty, zero divisor. Extend it by
Corollary~\ref{cor:canonical-endpoint} and shrink the disk so that
its divisor is smooth over the base. The associated cyclic cover
\[
 p:\mathcal Z\longrightarrow\mathcal X,\qquad
 p_*\OO_{\mathcal Z}
 =\bigoplus_{j=0}^{q-1}K_{\mathcal X/\Delta}^{\otimes(-j)}
\]
is smooth and proper over this disk.

Each connected component of $Z_0$ belongs to class $\mathcal C$.
Indeed, pull the cover back to a compact Kähler modification
$\mu:T\to X_0$. The pullback is a closed subspace of the projective
completion of $\mu^*K_{X_0}$, which is a compact Kähler manifold.
Resolving the components that dominate $T$ gives compact Kähler
modifications of the components of $Z_0$; every component of the
finite cover $Z_0\to X_0$ dominates $X_0$.
By Fujiki's degeneration theorem \cite[Proposition~1.6]{Fujiki1978},
the Fr\"olicher spectral sequence degenerates at $E_1$ on each
component of $Z_0$. In the smooth proper family
$\mathcal Z\to\Delta$, Betti numbers are constant and Hodge numbers
are upper semicontinuous. For each total degree, the Fr\"olicher
inequality and equality at $0$ force every Hodge number to be
constant near $0$. In particular, each $h^\ell(Z_t,\OO_{Z_t})$
is constant.

The cyclic-cover decomposition gives
\[
 h^\ell(Z_t,\OO_{Z_t})
 =\sum_{j=0}^{q-1}h^\ell(X_t,K_{X_t}^{\otimes(-j)}).
\]
Each summand is upper semicontinuous, and their sum is constant;
hence each summand is constant. Taking $j=m-1$ and
$\ell=\dim X_0-i$, Serre duality proves the fixed-degree assertion.

Apply Lemma~\ref{lem:uniform-higher} to
$\mathcal F_m=K_{\mathcal X/\Delta}^{\otimes m}$ and to the
relatively generated canonical power chosen at the beginning.
It gives one disk for all $m\ge1$ and all $i$, together with local
freeness and base change. Serre duality identifies degrees $m$
and $1-m$, giving constancy for $m\le0$ on the same disk.
Grauert's theorem gives the remaining direct-image assertions.
\end{proof}

\begin{corollary}[Nef central adjoints]\label{cor:nef-central-adjoint}
Assume the hypotheses of Theorem~\ref{thm:endpoint}, with $X$ compact
Kähler and $A=K_X+B$ nef in place of semiample.
Then Corollaries~\ref{cor:uniform-log-endpoint}
and~\ref{cor:relative-adjoint-algebra} apply.
If the family is smooth and $\mathcal B=0$, the conclusions of
Corollary~\ref{cor:smooth-semiample-higher} also hold.
\end{corollary}

\begin{proof}
Proposition~\ref{prop:decomposition-input} makes the nef adjoint
$A$ semiample. Apply the indicated corollaries.
\end{proof}

In particular, Theorem~\ref{thm:canonical-central} establishes the
constancy of plurigenera requested in
\cite[Question~1.6]{CRW2026}; it applies to the compact Kähler
central fiber in that question in every Kodaira dimension.

\subsection{Numerical consequences}

Hacon--Xie prove invariance of canonical pseudoeffectivity and
uniruledness for families with Kähler total space
\cite[Corollary~7.8]{HX2026}, resolving
\cite[Conjecture~1.4]{HLR2026}. The following corollary gives these
conclusions for families of compact Kähler canonical fibers.

\begin{corollary}[Numerical and birational invariants]
\label{cor:numerical-invariants}
Let $f:\mathcal X\to S$ be a proper flat morphism of complex
analytic spaces with $S$ connected and not necessarily reduced.
Suppose that all fibers are
connected $n$-dimensional compact Kähler spaces with canonical
singularities. Their Kodaira dimension
and canonical volume are constant. The loci
\[
 \{s\in S:K_{X_s}\text{ is pseudoeffective}\},\qquad
 \{s\in S:X_s\text{ is uniruled}\}
\]
are each either empty or equal to $S$.
\end{corollary}

\begin{proof}
The case $n=0$ is immediate; assume $n>0$.
Theorem~\ref{thm:canonical-central} and connectedness of $S$ make
each $P_m(X_s)$ constant. This sequence determines the Kodaira
dimension and the canonical volume
\[
 \vol(K_{X_s})
 :=\limsup_{m\to\infty}\frac{n!\,P_m(X_s)}{m^n}.
\]
For each fiber, Proposition~\ref{prop:decomposition-input} gives
\[
 K_{X_s}\text{ pseudoeffective}
 \quad\Longleftrightarrow\quad
 P_m(X_s)>0\text{ for some }m>0.
\]
For the forward implication, take a common multiple of the
semiample positive part and the effective rational negative part
in the good decomposition. A nonzero section of the resulting
pullback descends by the projection formula. The reverse implication
follows from the effective pluricanonical divisor.
Pseudoeffectivity is therefore constant.

Choose a compact Kähler resolution $W_s\to X_s$.
Canonical birational invariance gives
$\kappa(W_s)=\kappa(X_s)$. The same nonvanishing argument on $W_s$
identifies $\kappa(W_s)=-\infty$ with non-pseudoeffectivity of
$K_{W_s}$. Ou's criterion \cite[Theorem~1.1]{Ou2025} identifies
this condition with uniruledness of $W_s$, and uniruledness is
birationally invariant. It is consequently constant on the fibers
$X_s$ as well.
\end{proof}

\section{Finite logarithmic programs}
\label{sec:finite}

For logarithmic invariance, we use divisorial contractions and
flips between globally strongly $\Q$-factorial terminal pairs.
The scaling construction of \cite[Lemmas~3.6--3.7]{OAI034}
and Proposition~\ref{prop:termination-criterion} produce a finite
program starting on the prepared central pair.

\subsection{Ordinary logarithmic steps}

A normal space is \emph{globally strongly $\Q$-factorial} if every
global coherent rank-one reflexive sheaf has an invertible reflexive
power. In particular every global Weil divisor is $\Q$-Cartier.
This property holds on a smooth space. An \emph{ordinary negative
step} for $J=K_X+B$ means a projective divisorial contraction or a
projective small contraction and its flip, with normal contraction
base and connected fibers, such that $-J$ is relatively ample on the
contracting side and the transformed adjoint is relatively ample on
the positive side. Boundaries are pushed forward or strictly
transformed.

\begin{lemma}[Ordinary logarithmic scaling]
\label{lem:ordinary-scaling}\label{in:scaling}
Let $(X,B)$ be a globally strongly $\Q$-factorial compact Kähler klt
pair with effective rational boundary, and put $J=K_X+B$.
For a sufficiently large Kähler class $A$, there is an ordinary
$J$-negative program with scaling of $A$. Its
working spaces remain compact Kähler and globally strongly
$\Q$-factorial. Its maps are nonextracting and every positive truncation is finite. If $J$ is pseudoeffective, the
program either reaches a nef adjoint or has positive thresholds
tending to zero.

On a common resolution of a finite prefix, the unscaled comparison is
\begin{equation}\label{eq:ordinary-comparison}
 p^*J\sim_{\Q}q^*J_i+E_i,
 \qquad E_i\ge0,\quad q_*E_i=0.
\end{equation}
The coefficient is strictly positive at every original prime
contracted by the prefix, and a nontrivial ordinary negative step has
a nonzero comparison divisor. Discrepancies do not decrease.
The transforms of Bott--Chern classes are linear. At a scaling wall
the two scaled pullbacks agree.
\end{lemma}

\begin{proof}
These are the scaling and comparison statements of
\cite[Lemmas~3.6--3.7]{OAI034}, with the ordinary-step conventions
of Section~3.1 of that paper.
\end{proof}

\Needspace{13\baselineskip}
\subsection{The finite program}

\begin{proposition}\label{prop:finite}
Let $(X,B)$ be a smooth compact Kähler terminal pair with effective
rational SNC boundary. Assume that $J=K_X+B$ is pseudoeffective and
$N(J)\wedge B=0$. Then there is a finite sequence of ordinary
$J$-negative steps beginning at this $X$ and ending at a compact Kähler
terminal pair $(M,B_M)$ with $K_M+B_M$ semiample. No boundary component
is contracted. All working spaces are globally strongly
$\Q$-factorial, and their underlying spaces are terminal.
\end{proposition}

\begin{proof}
We apply Proposition~\ref{prop:termination-criterion} to the class
of globally strongly $\Q$-factorial compact Kähler terminal pairs
$(T,B_T)$ with effective rational boundary, pseudoeffective
$J_T=K_T+B_T$, and $N(J_T)\wedge B_T=0$.
Lemma~\ref{lem:ordinary-scaling} gives the scaling programs and comparisons;
pseudoeffectivity persists by pushing forward a positive current
representing the original adjoint on a common resolution.

We verify that the working pairs remain in this class. Strict
comparison and exceptional translation place each contracted prime
in the negative part of the current adjoint. Thus no boundary prime
is contracted. On a common resolution, the same comparison shows
that minimal multiplicities along surviving boundary primes remain
zero. An exceptional valuation over the initial model has positive
discrepancy, which cannot decrease. A prime of the initial model
contracted later has boundary coefficient zero, and strict comparison
raises its discrepancy from zero to a positive value. These are all
exceptional valuations over a working model, so every working pair
is terminal. Global strong $\Q$-factoriality makes its boundary
$\Q$-Cartier; removing this effective boundary increases exceptional
discrepancies, hence its underlying space is terminal as well.

Condition~\textup{(i)} of
Proposition~\ref{prop:termination-criterion} is therefore satisfied,
and Proposition~\ref{prop:interval-input} supplies condition~\textup{(ii)}.
The criterion gives a finite ordinary program starting on $X$.
Proposition~\ref{prop:decomposition-input} makes its nef endpoint
semiample.
\end{proof}


\section{Preparation and deformation of logarithmic programs}
\label{sec:relative-reduction}

We prepare the central pair in two steps: crepant stratum blowups
make it terminal, and subtraction of the fixed boundary separates
the remaining boundary from the divisorial negative part. Both
steps preserve the prescribed central section space. The resulting
finite logarithmic program then deforms over the disk.

Throughout this section the base is a disk and the fibers are
connected. After shrinking, we split every boundary component and
closed boundary stratum according to the connected components of
its fibers. These are proper smooth families, so their sets of
connected components are locally constant. Each resulting relative
boundary component therefore restricts to one central prime divisor.

\Needspace{19\baselineskip}
\subsection{Simultaneous terminal preparation}

The following terminal preparation refines the crepant
stratum-blowup construction of \cite[Section~4.2]{LW2026}.
Keeping the coefficients in $m^{-1}\Z$ preserves the integral
adjoint in the prescribed degree.

\begin{lemma}\label{lem:terminal-prep}
Let $f:\cX\to\Delta$ be a smooth proper holomorphic family with
compact Kähler fibers and relative SNC boundary
$B=\sum b_iD_i$, where $0\le b_i<1$ and $mb_i\in\Z$ for a fixed
positive integer $m$. After shrinking the disk there is a finite
composition of blowups of smooth relative codimension-two boundary
strata,
\[
 \tau:\widetilde{\cX}\longrightarrow\cX,
\]
with a relative SNC boundary $\widetilde B$ such that
\begin{equation}\label{eq:terminal-preparation}
 K_{\widetilde{\cX}/\Delta}+\widetilde B
 =\tau^*(K_{\cX/\Delta}+B),\qquad
 0\le\widetilde B<1,
 \qquad m\widetilde B\text{ is integral}.
\end{equation}
The new family is smooth with compact Kähler fibers. Its total pair
and every fiber pair are terminal. Pullback identifies the degree-$m$
adjoint sections, on the total spaces and on every fiber.
\end{lemma}

\begin{proof}
\emph{The blowups.}
Whenever two intersecting boundary components have coefficients
$b_i,b_j$ with $b_i+b_j\ge1$, blow up a connected component of their
intersection and give the new exceptional divisor coefficient
$b_i+b_j-1$. The center is smooth of relative codimension two, and
the blowup formula gives \eqref{eq:terminal-preparation} at this
step. The new coefficient lies in $[0,1)$ and is an integral multiple
of $1/m$. The family stays smooth and the total boundary support,
including any zero-coefficient exceptional components temporarily
retained, stays relative SNC.

\smallskip\noindent\emph{Termination.}
Consider the finite cone complex defined by the
original relative SNC strata after shrinking around the compact
central fiber. Each closed intersection has finitely many components,
proper and smooth over the disk. In an original cone use its boundary generators as
the lattice basis and put $d_i=1-b_i\ge1/m$. The blowups above are
star subdivisions along sums of two generators of a regular cone.
Thus every inserted ray is primitive, has nonnegative integral
coordinates $v=(v_i)$ in an original cone, and has log discrepancy
\[
 A(v)=\sum_i v_id_i.
\]
The coefficient assigned to that ray is $1-A(v)$. It is inserted only
when this coefficient is nonnegative, so
\[
 \sum_i v_i\le m.
\]
There are finitely many such lattice vectors in each original cone.
A subdivision inserts a ray not already present in the cone complex,
and subsequent subdivisions retain every inserted ray. Keeping the
finitely many original stratum labels distinguishes rays occurring
over different connected components. The process therefore terminates.

\smallskip\noindent\emph{Terminality and section comparison.}
At termination, any two intersecting positive boundary coefficients
have sum less than one. Let $v$ be an exceptional divisorial
valuation over a
smooth model in question, and take a general smooth point of its
center, of codimension $c\ge2$. Choose local coordinates such that
$z_1,\ldots,z_c$ vanish on the center and the boundary components
containing it are $z_1=0,\ldots,z_s=0$. Write $v_i=v(z_i)\ge1$ for
$i\le c$. Pulling back the coordinate volume form to a model carrying
$v$ shows
\[
 A_X(v)\ge\sum_{i=1}^c v_i:
\]
in its exterior product at most one differential can lower the order
by contributing the differential of a local equation of the
exceptional prime. Consequently
\begin{equation}\label{eq:snc-terminal-valuation}
 a(v,X,B)\ge
 \sum_{i=1}^s(1-b_i)v_i+\sum_{i=s+1}^c v_i-1>0.
\end{equation}
If $s\ge2$, two of the first coefficients already sum to more than
one. If $s=1$, there is at least one further positive summand; if
$s=0$, the sum is at least two. This proves the strict inequality in
every case. The same argument applies to the total smooth space and
to every fiber; compare \cite[Corollary~2.31]{KM1998}.

Each fiberwise blowup is projective over a compact Kähler manifold
and is again Kähler. Finally, multiplying the crepant identity by $m$ gives an
isomorphism of the integral adjoint line bundles. The equality
$\tau_*\OO_{\widetilde{\cX}}=\OO_{\cX}$ and the projection formula
give the section identification on the total spaces. The same
argument applies to each fiberwise blowup.
\end{proof}

\subsection{The fixed boundary in a prescribed degree}

We next apply the degreewise fixed-part argument of
\cite[proof of Theorem~4.2]{HMX2018}; compare
\cite[Section~4.3]{LW2026}.

\begin{lemma}\label{lem:fixed-boundary}
Let $(\cX,B)\to\Delta$ be a prepared family as in
Lemma~\ref{lem:terminal-prep}. Suppose
\[
 H^0\bigl(X_0,\OO_{X_0}(m(K_{X_0}+B_0))\bigr)\ne0.
\]
There is a relative boundary $B_*\le B$, with $mB_*$ integral and
with terminal total and fiber pairs, such that
\[
 J_*=K_{X_0}+B_{*,0}\text{ is pseudoeffective},
 \qquad N(J_*)\wedge B_{*,0}=0.
\]
For the relative effective integral divisor $F=m(B-B_*)$,
multiplication by its canonical section identifies the central
section spaces:
\begin{equation}\label{eq:fixed-boundary-sections}
 H^0\bigl(X_0,m(K_{X_0}+B_0)\bigr)
 =s_{F_0}\,H^0\bigl(X_0,m(K_{X_0}+B_{*,0})\bigr).
\end{equation}
\end{lemma}

\begin{proof}
Write $e_i$ for the fixed multiplicity along $D_{i,0}$ of the
complete central linear system, namely the least zero order there
among its nonzero sections. Set
\[
 q_i=\min\{mb_i,e_i\},\qquad
 F=\sum_iq_iD_i,\qquad B_*=B-F/m.
\]
Both $q_i$ and $e_i$ are nonnegative integers. Every central section
is divisible by $s_{F_0}$, and division produces a holomorphic section
of the smaller line bundle. Conversely multiplication by $s_{F_0}$
takes each section of that smaller line into the original complete
system. This proves \eqref{eq:fixed-boundary-sections}.

The new fixed multiplicity along $D_{i,0}$ is exactly $e_i-q_i$.
If its remaining boundary coefficient $b_i-q_i/m$ is positive, then
$q_i=e_i$. Some section of the smaller adjoint therefore has zero order
along this prime. Its zero divisor divided by $m$ represents a
positive current in $\{J_*\}$ with zero generic Lelong number there.
It follows that $\nu_{D_{i,0}}(J_*)=0$. This proves the stated
negative-part condition. The same nonzero section proves
pseudoeffectivity of $J_*$. Lowering an effective boundary increases
discrepancies, and all coefficients retain denominator dividing $m$.
\end{proof}

The relative divisor $F$ records the central fixed multiplicities.
Its section $s_F$ restores the original boundary after extension
for $B_*$.

\Needspace{27\baselineskip}
\subsection{Deformation of the prepared program}

\begin{proposition}\label{prop:relative-program}
Apply Proposition~\ref{prop:finite} to the central pair obtained from
Lemmas~\ref{lem:terminal-prep} and~\ref{lem:fixed-boundary}.
After shrinking the disk, its finite ordinary program extends to
proper relative bimeromorphic diagrams beginning at the prepared
family $(\cX,B_*)$. Their endpoint
$g:(\cY,\cB)\to\Delta$ has the following properties.
\begin{enumerate}[label=\textup{(\roman*)},leftmargin=*]
\item $\cY$ is normal and $g$ is proper and flat. Its central pair
$(Y_0,\cB_0)$ is the prescribed compact Kähler terminal endpoint.
In particular $K_{Y_0}+\cB_0$ is semiample.
\item The relative log adjoint
$\cA=K_{\cY/\Delta}+\cB$ is $\Q$-Cartier. Both the total pair
$(\cY,\cB)$ and the underlying total space are canonical near $Y_0$.
\item $\cB$ is the horizontal strict transform of $B_*$, its
coefficients lie in $[0,1)$, and $m\cB$ and $m\cB_0$ are integral.
Adjunction identifies $\cA|_{Y_0}$ with
$K_{Y_0}+\cB_0$ as holomorphic $\Q$-line bundles.
\item A degree-$m$ section on the endpoint total space transfers,
under the meromorphic pluricanonical identification, to a section
on the prepared family with boundary $B_*$. This transfer commutes
with central restriction and with the analogous identification of
central section spaces.
\end{enumerate}
\end{proposition}

\begin{proof}
\emph{Construction of the family.}
We apply \cite[Theorem~7.7]{HX2026} with generalized nef datum
zero. The prepared total space is smooth, the boundary is rational
and horizontal, and the log adjoint is $\Q$-Cartier.
After choosing a relatively compact smaller disk, the morphism is
flat and proper with connected fibers. The central pair is canonical,
has empty boundary floor, and satisfies
$N(K_{X_0}+B_{*,0})\wedge B_{*,0}=0$ by
Lemma~\ref{lem:fixed-boundary}. Thus the central hypotheses of that theorem hold.

Hacon--Xie extend the prescribed central ordinary steps. Their proof
uses the construction of \cite[Proposition~4.16]{HLR2026}, replacing
the latter's projective-central-fiber model input by their Kähler
model theorem. In that construction, contraction deformation gives
the contraction diagrams; in the small case, the local relative log
canonical models glue by uniqueness to give the flip. The construction
identifies the central contraction, central flip, or central
divisorial target with the prescribed one. It therefore gives proper relative diagrams with the prescribed
normal central spaces. Since the program is finite, one shrinking
of the disk suffices.
The total spaces are normal by the construction. On each total space dominating the disk, the parameter is a
nonzerodivisor. Torsion-freeness over the local discrete valuation
ring of the base therefore gives flatness. This proves the family assertions in (i).

\smallskip\noindent\emph{Boundary, singularities, and adjunction.}
A central divisorial contraction is an isomorphism at the generic
point of every prime of its normal target; the two small contractions
of a flip have the same property. Lemma~\ref{lem:open-isomorphism}
extends each of these isomorphism loci to the total spaces. Applying
it successively to the finite diagrams shows that their common
isomorphism locus meets every endpoint central prime.

At these generic points, strict transform preserves the boundary
coefficients and gives multiplicity-one central restrictions.
Every endpoint central boundary prime is thus the transform of an
original central boundary prime, and Proposition~\ref{prop:finite}
preserves all such primes. The total boundary remains horizontal
and its coefficients remain in $m^{-1}\Z\cap[0,1)$.

The canonical-pair deformation in
\cite[Proposition~4.16]{HLR2026}, used
by \cite[Theorem~7.7]{HX2026}, gives canonical total ordinary pairs
and an $\R$-Cartier log adjoint near their central fibers. This applies to a normal Cartier central fiber with canonical
restricted pair and boundary coefficients less than one; compare \cite[Proposition~3.5]{dFH2011} and the analytic
Cartier-extension arguments in \cite[Section~12.1]{KM1992}.
In a local meromorphic canonical frame, the log adjoint is a
rational Weil divisor. Since it is $\R$-Cartier, its local divisor
class vanishes after tensoring with $\R$, hence also after
tensoring with $\Q$. Thus the adjoint is locally $\Q$-Cartier. Compactness gives
a common Cartier multiple near the central fiber.

Proposition~\ref{prop:finite} also makes each underlying central
working space terminal. Deformation of canonical singularities
therefore gives a $\Q$-Cartier relative canonical divisor and
canonical underlying total space near the central fiber;
see \cite[Lemma~10 and Remark~10.3]{Kollar2021}, based on
\cite[Theorem~VI.5.2]{Nakayama2004}. In particular the total space
is klt.

At every generic central prime, the local isomorphism identifies
ordinary adjunction and the multiplicity-one boundary restriction. Taking a common Cartier multiple extends that
identification over the normal central fiber. This is the
$\Q$-linear identification claimed in (iii). Lemma~\ref{lem:restrictions} consequently identifies
$\OO_{\cY}(m\cA)/t\OO_{\cY}(m\cA)$ with
$\OO_{Y_0}(m(K_{Y_0}+\cB_0))$. This proves (ii) and (iii).

\smallskip\noindent\emph{Transfer of sections.}
Apply Lemma~\ref{lem:birational} to the canonical total
pairs at the start and end of the diagrams. Their boundaries agree
on common nonexceptional primes, and the prescribed degree is
admissible on both sides. This gives the section transfer in (iv).
The same lemma applies to the canonical central pairs. The two
transfers commute with restriction on the common dense smooth open
set, where they are ordinary pullbacks of relative pluriforms.
The quotient identification above and reflexivity extend this
commutativity to the whole central fiber.
\end{proof}

\section{Proof of logarithmic invariance}\label{sec:main-proof}

\begin{proof}[Proof of Theorem~\ref{thm:main}]
\emph{Extension over a disk.}

We first prove local extension over a disk. After shrinking, the
connected components of a smooth proper family split, so we work
with connected fibers. The zero-dimensional case is immediate.

Fix $m\ge1$ and set
\begin{equation}\label{eq:degree-rounding}
 B^{(m)}=\frac1m\lfloor mB\rfloor.
\end{equation}
This is a rational relative SNC boundary with coefficients in $[0,1)$,
and
\begin{equation}\label{eq:rounded-identity}
 \mathcal L_m=\OO_{\mathcal X}
       \bigl(m(K_{\mathcal X/\Delta}+B^{(m)})\bigr).
\end{equation}
For $m=1$, the boundary $B^{(1)}$ is zero and
Theorem~\ref{thm:degreeone} gives extension. For $m\ge2$,
it suffices to extend an arbitrary nonzero section
$u_0\in H^0(X_0,\mathcal L_m|_{X_0})$.

Apply Lemma~\ref{lem:terminal-prep} to obtain a crepant relative
preparation
\[
 \tau:(\widetilde{\mathcal X},\widetilde B)
       \longrightarrow(\mathcal X,B^{(m)}).
\]
Lemma~\ref{lem:fixed-boundary} then gives a boundary
$\widetilde B_*\le\widetilde B$ and an effective relative integral
divisor $F=m(\widetilde B-\widetilde B_*)$ such that
\[
 J_*=K_{\widetilde X_0}+\widetilde B_{*,0}
 \quad\text{is pseudoeffective},\qquad
 N(J_*)\wedge\widetilde B_{*,0}=0.
\]
The boundary coefficients still have denominator dividing $m$.
Proposition~\ref{prop:finite} supplies a finite central program with
terminal semiample endpoint $(Y_0,B_{Y_0})$.
Proposition~\ref{prop:relative-program} extends it to a proper flat
family $g:(\mathcal Y,\mathcal B)\to\Delta$ and verifies the
hypotheses of Theorem~\ref{thm:endpoint}, including adjunction and
the integrality of $m\mathcal B$ and $mB_{Y_0}$.

The crepant preparation, fixed-boundary identity, and
birational comparison give
\[
 \frac{\tau_0^*u_0}{s_{F_0}}
 \in H^0(\widetilde X_0,mJ_*)
 \simeq H^0(Y_0,m(K_{Y_0}+B_{Y_0})).
\]
Theorem~\ref{thm:endpoint} extends the corresponding endpoint
section in degree $m$. Transfer it back by
Proposition~\ref{prop:relative-program}(iv), obtaining
\[
 v\in H^0\bigl(\widetilde{\mathcal X},
       m(K_{\widetilde{\mathcal X}/\Delta}+\widetilde B_*)\bigr).
\]
Multiplication by $s_F$ restores the boundary $\widetilde B$.
The resulting section descends through $\tau$ by the integral
crepant identity and the projection formula. Its central restriction
is $u_0$. Thus every central section extends in its prescribed degree.

The line bundle $\mathcal L_m$ is flat over the disk, so
Lemma~\ref{lem:extension-constancy} gives local constancy of
$h^0(X_t,\mathcal L_m|_{X_t})$. Repeating the argument at every
point proves constancy on a connected disk.

\smallskip\noindent\emph{Smooth bases and relative consequences.}

For a connected complex manifold $S$, restrict to coordinate disks
in a coordinate polydisk. Each restricted family remains proper and
smooth, with compact Kähler fibers and relative SNC boundary, and
\eqref{eq:rounded-identity} restricts to its degree-$m$ line bundle.
The disk result makes the dimension constant along each such disk.
Any two points of a smaller coordinate polydisk can be joined by
finitely many such segments. The dimension is therefore locally
constant, and connectedness makes it constant on $S$.

Grauert's theorem \cite{Grauert1960} now gives local freeness and,
for every holomorphic map $a:T\to S$, the base-change isomorphism
\begin{equation}\label{eq:all-base-change}
 a^*(f_*\mathcal L_m)
 \xrightarrow{\ \sim\ }
 (f_T)_*(p_{\mathcal X}^*\mathcal L_m),
 \qquad
 p_{\mathcal X}:\mathcal X\times_ST\to\mathcal X.
\end{equation}
In particular the fiber of $f_*\mathcal L_m$ at $s$ is
$H^0(X_s,\mathcal L_m|_{X_s})$. A local frame therefore extends
each fiber section over a neighborhood of $s$.
If $S$ is Stein, Cartan's theorem~A makes global sections of
$f_*\mathcal L_m$ surject onto each fiber. Since
$H^0(S,f_*\mathcal L_m)=H^0(\mathcal X,\mathcal L_m)$,
the extension can then be taken over all of $\mathcal X$.
\end{proof}

\section{Twistor families and the Kähler hypotheses}
\label{sec:example}

The twistor construction gives smooth families whose fibers are compact
Kähler and whose total spaces are non-Kähler over every open subset
of the base. The period description makes this local property explicit.

\begin{proposition}\label{prop:twistor-local}
Let $q:Z\to\PP^1$ be the twistor family of a hyperkähler K3 surface.
For every nonempty connected open subset $U\subset\PP^1$, the
restriction $Z_U\to U$ is a proper holomorphic submersion with
compact Kähler fibers, whereas $Z_U$ is not Kähler.
\end{proposition}

\begin{proof}
Let $M$ be the underlying smooth four-manifold, with hyperkähler
metric $g$ and quaternionic complex structures $I,J,K$. Put
\[
 \omega_I=g(I\,\cdot,\cdot),\quad
 \omega_J=g(J\,\cdot,\cdot),\quad
 \omega_K=g(K\,\cdot,\cdot).
\]
The twistor construction gives the differentiable product
$M\times\PP^1$ an integrable complex structure for which the second
projection is a holomorphic submersion. The fiber indexed by
$L=aI+bJ+cK$, $a^2+b^2+c^2=1$, has complex structure $L$ and Kähler
form
\[
 \omega_L=a\omega_I+b\omega_J+c\omega_K.
\]
We use the standard construction and its period description from
\cite[Chapter~7, Section~3.2]{Huybrechts2016}; see also
\cite[Section~1.3]{Huybrechts2023}. Compactness of $M$ gives
properness, including after restriction to $U$.

Under the differentiable identification of the fiber cohomologies,
let
\[
 W=\operatorname{span}_{\R}
   \{[\omega_I],[\omega_J],[\omega_K]\}\subset H^2(M,\R).
\]
The intersection form is positive definite on $W$. Indeed
quaternionic linear algebra gives
\[
 \omega_I\wedge\omega_J
 =\omega_J\wedge\omega_K
 =\omega_K\wedge\omega_I=0,
 \qquad
 \omega_I^2=\omega_J^2=\omega_K^2=2\vol_g.
\]
For the complex structure $L$, the real $(2,0)+(0,2)$ plane is
\begin{equation}\label{eq:twistor-period-plane}
 P_L=W\cap[\omega_L]^\perp.
\end{equation}
For example, complete $L$ to an oriented orthonormal quaternionic
triple $L,L',L''$; the holomorphic symplectic form is
$\omega_{L'}+\sqrt{-1}\omega_{L''}$, proving
\eqref{eq:twistor-period-plane}.

Suppose that $\Omega$ were a Kähler form on $Z_U$. Since it is
closed and the fiber inclusions in $M\times U$ are homotopic along
paths, their restrictions have a fixed de Rham class
$\alpha\in H^2(M,\R)$. Each restriction is of type $(1,1)$ for its
fiber. Hodge orthogonality and \eqref{eq:twistor-period-plane}
therefore imply that the $W$-projection of $\alpha$ lies in
$\R[\omega_L]$ for every $L\in U$. The open set $U$ contains two
distinct nonantipodal points of the sphere, whose two lines have
zero intersection. Thus this projection is zero, and
\[
 \int_M\alpha\wedge[\omega_L]=0\qquad(L\in U).
\]
On the other hand, $\Omega|_{(M,L)}$ and $\omega_L$ are Kähler forms
on a compact complex surface. Their wedge is a strictly positive
volume form, so this integral is positive. The contradiction proves
the proposition.
\end{proof}

\begin{example}[A nonzero relative boundary]\label{ex:twistor-boundary}
Let $U$ be a disk in the twistor sphere, let $C$ be a smooth projective
curve of genus $g\ge2$, and choose $r\ge1$ distinct points $p_1,\ldots,p_r$
and rational coefficients $0<b_i<1$. The family
\[
 \cX=Z_U\times C\longrightarrow U,
 \qquad D_i=Z_U\times\{p_i\},\qquad B=\sum_i b_iD_i
\]
has compact Kähler fibers and a nonzero relative SNC boundary with
disjoint components. Its total space is not Kähler over any smaller
nonempty disk: restriction of a hypothetical Kähler form to
$Z_U\times\{p\}$ would contradict Proposition~\ref{prop:twistor-local}.

These examples have nonzero log pluricanonical spaces in every degree.
The canonical line of every K3 fiber is trivial, and its pluricanonical
space is one-dimensional. The product formula for sections therefore
gives
\[
 h^0\bigl(X_t,mK_{X_t}+\lfloor mB_t\rfloor\bigr)
 =h^0\left(C,mK_C+\sum_i\lfloor mb_i\rfloor p_i\right).
\]
For $m=1$ the right side is $g$. For $m\ge2$, Serre duality gives
zero first cohomology, since
$(1-m)K_C-\sum_i\lfloor mb_i\rfloor p_i$ has negative degree.
Riemann--Roch gives the explicit value
\[
 (2m-1)(g-1)+\sum_i\lfloor mb_i\rfloor.
\]
The formula computes the logarithmic plurigenera of this family in
every degree and exhibits their constancy on $U$.
\end{example}

\appendix
\section{Abundance and nef comparison}\label{app:dependencies}
\label{app:canonical-dependencies}

The termination arguments use two consequences of
\cite{OAI034}. Its abundance theorem
\cite[Theorem~2.13]{OAI034} applies by logarithmic subadditivity
\cite[Corollary~6.2]{OAI033}. We state the resulting decomposition
and nef comparison in the form used here; all decompositions are
$\Q$-linear equivalences of holomorphic $\Q$-line bundles.

\begin{proposition}[Good divisorial decomposition]
\label{prop:decomposition-input}\label{in:good}
Let $(X,B)$ be a normal compact Kähler log canonical pair, with
$B$ an effective rational divisor and $J=K_X+B$ a pseudoeffective
$\mathbb Q$-Cartier adjoint. There is a smooth compact Kähler
modification $\mu:W\to X$ such that
\[
 \mu^*J\sim_{\mathbb Q}P+N(\mu^*J),
 \qquad P\text{ semiample},
\]
and $N(\mu^*J)$ is rational. On every higher smooth model
$a:W'\to W$, the negative part is $a^*N(\mu^*J)$.
\end{proposition}

\begin{proof}
Apply \cite[Theorem~2.13]{OAI034}, with the higher-pullback and
normal-pair formulations of \cite[Lemma~2.4 and Proposition~2.7]{OAI034}.
\end{proof}

In particular, a nef adjoint in Proposition~\ref{prop:decomposition-input}
is semiample: its negative part vanishes after pullback, and generation
of a Cartier multiple descends by normality and the projection formula.

\begin{proposition}[Nef comparison near zero]
\label{prop:interval-input}\label{in:interval}
Let $(T,B_T)$ be a normal compact Kähler klt pair with effective
rational boundary, and put $J_T=K_T+B_T$. Suppose that $J_T$ is
$\mathbb Q$-Cartier and pseudoeffective, and that $N(J_T)=0$ on
the prime divisors of $T$. For every Kähler class $A$ there are
a normal compact Kähler space $V$, a small bimeromorphic map
$T\dashrightarrow V$, transformed classes $\{J_V\},A_V$, and
$\delta>0$ such that
\[
 \{J_V\}+sA_V\quad\text{is nef for }0\le s\le\delta.
\]
On a common smooth resolution, the difference of the original and
transformed pullbacks is represented by a real divisor exceptional
over both models, depending affinely on $s$.
\end{proposition}

\begin{proof}
This is the empty-floor case of
\cite[Lemmas~3.11--3.12]{OAI034}. Those lemmas apply to normal dlt pairs. Here the pair is klt,
so their auxiliary divisor supported on $\lfloor B_T\rfloor$
is zero.
Their dimensional hypothesis is satisfied with
$n=\dim T+1$, since \cite[Theorem~2.13]{OAI034} supplies good
decomposition in every dimension.
\end{proof}

The condition $N(J_T)=0$ concerns the prime divisors of $T$,
which is exactly the vanishing obtained in
Proposition~\ref{prop:termination-criterion} before restarting
the scaling program.

\section*{AI Disclosure}
All mathematical derivations and proofs were produced by AI. The human
author formulated the problem, provided relevant literature, and guided
the research direction. AI also assisted in the selection of some of the
problems.
\begin{thebibliography}{99}
\small
\setlength{\itemsep}{4pt}
\setlength{\parsep}{0pt}
\setlength{\parskip}{0pt}
\raggedright
\interlinepenalty=10000

\bibitem{BM1997}
E.~Bierstone and P.~D.~Milman,
\emph{Canonical desingularization in characteristic zero by
blowing up the maximum strata of a local invariant},
Invent. Math. \textbf{128} (1997), no.~2, 207--302.
\href{https://doi.org/10.1007/s002220050141}{doi:10.1007/s002220050141}.

\bibitem{Bingener1983}
J.~Bingener,
\emph{On deformations of Kähler spaces. I},
Math. Z. \textbf{182} (1983), no.~4, 505--535.
\href{https://doi.org/10.1007/BF01215480}{doi:10.1007/BF01215480}.

\bibitem{BCHM2010}
C.~Birkar, P.~Cascini, C.~D.~Hacon, and J.~McKernan,
\emph{Existence of minimal \mbox{models} for varieties of log general type},
J. Amer. Math. Soc. \textbf{23} (2010), no.~2, 405--468.
\href{https://doi.org/10.1090/S0894-0347-09-00649-3}{doi:10.1090/S0894-0347-09-00649-3}.

\bibitem{Boucksom}
S.~Boucksom,
\emph{Divisorial Zariski decompositions on compact complex manifolds},
Ann. Sci. École Norm. Sup. (4) \textbf{37} (2004), no.~1, 45--76.
\href{https://doi.org/10.1016/j.ansens.2003.04.002}{doi:10.1016/j.ansens.2003.04.002}.

\bibitem{CaoPaun2023}
J.~Cao and M.~Păun,
\emph{Infinitesimal extension of pluricanonical forms},
preprint, 2023.
\href{https://arxiv.org/abs/2012.05063v5}{arXiv:2012.05063v5}.

\bibitem{CRW2026}
J.~Chen, S.~Rao, and K.~Wang,
\emph{Matsumura's extension problem for pluricanonical forms in
Kähler families I: the smooth and essentially Moishezon cases},
preprint, 30 September 2026.
\href{https://arxiv.org/abs/2609.40040v1}{arXiv:2609.40040v1}.

\bibitem{dFH2011}
T.~de Fernex and C.~D.~Hacon,
\emph{Deformations of canonical pairs and Fano varieties},
J. Reine Angew. Math. \textbf{651} (2011), 97--126.
\href{https://doi.org/10.1515/CRELLE.2011.010}{doi:10.1515/CRELLE.2011.010};
\href{https://arxiv.org/abs/0901.0389v3}{\mbox{arXiv:0901.0389v3}}.

\bibitem{Deligne1971}
P.~Deligne,
\emph{Théorie de Hodge, II},
Publ. Math. IHÉS \textbf{40} (1971), 5--57.
\href{https://doi.org/10.1007/BF02684692}{doi:10.1007/BF02684692}.

\bibitem{DNWZ2026}
F.~Deng, J.~Ning, Z.~Wang, and X.~Zhou,
\emph{Invariance of plurigenera for Kähler families with nef canonical bundles},
preprint, 29 September 2026.
\href{https://arxiv.org/abs/2609.37804v1}{arXiv:2609.37804v1}.

\bibitem{EV1992}
H.~Esnault and E.~Viehweg,
\emph{Lectures on Vanishing Theorems},
DMV Seminar, vol.~20, Birkhäuser Verlag, Basel, 1992.
\href{https://doi.org/10.1007/978-3-0348-8600-0}{doi:10.1007/978-3-0348-8600-0}.

\bibitem{FLR2026}
J.~Feng, Y.~Li, and S.~Rao,
\emph{Aspects of deformation invariance of plurigenera of Kähler varieties},
preprint, 13 October 2026.
\href{https://arxiv.org/abs/2610.17817v1}{arXiv:2610.17817v1}.

\bibitem{Frisch1967}
J.~Frisch,
\emph{Points de platitude d'un morphisme d'espaces analytiques complexes},
Invent. Math. \textbf{4} (1967), 118--138.
\href{https://doi.org/10.1007/BF01425245}{doi:10.1007/BF01425245}.

\bibitem{Fujiki1978}
A.~Fujiki,
\emph{On automorphism groups of compact Kähler manifolds},
Invent. Math. \textbf{44} (1978), no.~3, 225--258.
\href{https://doi.org/10.1007/BF01403162}{doi:10.1007/BF01403162}.

\bibitem{Fujino2022}
O.~Fujino,
\emph{Minimal model program for projective morphisms between complex
analytic spaces}, preprint, 2022.
\href{https://arxiv.org/abs/2201.11315v1}{arXiv:2201.11315v1}.

\bibitem{Grauert1960}
H.~Grauert,
\emph{Ein Theorem der analytischen Garbentheorie und die Modulräume
komplexer Strukturen},
Publ. Math. IHÉS \textbf{5} (1960), 5--64.
\href{https://doi.org/10.1007/BF02684746}{doi:10.1007/BF02684746}.
Berichtigung, Publ. Math. IHÉS \textbf{16} (1963), 35--36.
\href{https://doi.org/10.1007/BF02684291}{doi:10.1007/BF02684291}.

\bibitem{HLR2026}
C.~D.~Hacon, Y.~Li, and S.~Rao,
\emph{On pseudo-effectivity and volumes of adjoint classes in Kähler
families with projective central fiber},
preprint, 2026, revised 5 March 2026.
\href{https://arxiv.org/abs/2602.04158v2}{arXiv:2602.04158v2}.

\bibitem{HLX2026}
C.~D.~Hacon, Y.~Li, and L.~Xie,
\emph{Fujiki class $\mathcal C$ varieties and a Kähler criterion},
preprint, 2026.
\href{https://arxiv.org/abs/2608.20588v1}{arXiv:2608.20588v1}.

\bibitem{HMX2018}
C.~D.~Hacon, J.~McKernan, and C.~Xu,
\emph{Boundedness of moduli of varieties of general type},
J. Eur. Math. Soc. \textbf{20} (2018), no.~4, 865--901.
\href{https://doi.org/10.4171/JEMS/778}{doi:10.4171/JEMS/778}.

\bibitem{HP2024}
C.~D.~Hacon and M.~Păun,
\emph{On the canonical bundle formula and adjunction for generalized
Kähler pairs}, preprint, 2024.
\href{https://arxiv.org/abs/2404.12007v1}{arXiv:2404.12007v1}.

\bibitem{HX2026}
C.~D.~Hacon and L.~Xie,
\emph{On the Kähler MMP and the transcendental base-point-free theorem},
preprint, 27 July 2026.
\href{https://arxiv.org/abs/2607.24986v1}{arXiv:2607.24986v1}.

\bibitem{HH2020}
K.~Hashizume and Z.-Y.~Hu,
\emph{On minimal model theory for log abundant lc pairs},
J. Reine Angew. Math. \textbf{767} (2020), 109--159.
\href{https://doi.org/10.1515/crelle-2019-0032}{doi:10.1515/crelle-2019-0032};
\href{https://arxiv.org/abs/1906.00769v4}{arXiv:1906.00769v4}.

\bibitem{Hironaka1975}
H.~Hironaka,
\emph{Flattening theorem in complex-analytic geometry},
Amer. J. Math. \textbf{97} (1975), no.~2, 503--547.
\href{https://doi.org/10.2307/2373721}{doi:10.2307/2373721}.

\bibitem{Huybrechts2016}
D.~Huybrechts,
\emph{Lectures on K3 Surfaces},
Cambridge Studies in Advanced Mathematics, vol.~158,
Cambridge University Press, Cambridge, 2016.
\href{https://doi.org/10.1017/CBO9781316594193}{doi:10.1017/CBO9781316594193}.

\bibitem{Huybrechts2023}
D.~Huybrechts,
\emph{Brilliant families of K3 surfaces: Twistor spaces, Brauer groups,
and Noether--Lefschetz loci},
Ann. Fac. Sci. Toulouse Math. (6) \textbf{32} (2023), no.~2, 397--421.
\href{https://doi.org/10.5802/afst.1741}{doi:10.5802/afst.1741}.

\bibitem{Iitaka1970}
S.~Iitaka,
\emph{Deformations of compact complex surfaces II},
J. Math. Soc. Japan \textbf{22} (1970), no.~2, 247--261.
\href{https://doi.org/10.2969/jmsj/02220247}{doi:10.2969/jmsj/02220247}.

\bibitem{Kawamata1999}
Y.~Kawamata,
\emph{Deformations of canonical singularities},
J. Amer. Math. Soc. \textbf{12} (1999), no.~1, 85--92.
\href{https://doi.org/10.1090/S0894-0347-99-00285-4}{doi:10.1090/S0894-0347-99-00285-4}.

\bibitem{Kirschner2012}
T.~Kirschner,
\emph{Irreducible Symplectic Complex Spaces},
Ph.D. thesis, Universität Bayreuth, 2012.
\href{https://epub.uni-bayreuth.de/id/eprint/209/1/diss-kirschner-bibprint.pdf}{Dissertation}.

\bibitem{Kollar2021}
J.~Kollár,
\emph{Deformations of varieties of general type},
Milan J. Math. \textbf{89} (2021), no.~2, 345--354.
\href{https://doi.org/10.1007/s00032-021-00339-6}{doi:10.1007/s00032-021-00339-6};
\href{https://arxiv.org/abs/2101.10986v3}{arXiv:2101.10986v3}.

\bibitem{KM1992}
J.~Kollár and S.~Mori,
\emph{Classification of three-dimensional flips},
J. Amer. Math. Soc. \textbf{5} (1992), no.~3, 533--703.
\href{https://doi.org/10.1090/S0894-0347-1992-1149195-9}{doi:10.1090/S0894-0347-1992-1149195-9}.
See also S.~Mori, \emph{Errata to “Classification of three-dimensional flips”},
J. Amer. Math. Soc. \textbf{20} (2007), no.~1, 269--271.
\href{https://doi.org/10.1090/S0894-0347-06-00543-1}{doi:10.1090/S0894-0347-06-00543-1}.

\bibitem{KM1998}
J.~Kollár and S.~Mori,
\emph{Birational Geometry of Algebraic Varieties},
Cambridge Tracts in Mathematics, vol.~134,
Cambridge University Press, Cambridge, 1998.
\href{https://doi.org/10.1017/CBO9780511662560}{doi:10.1017/CBO9780511662560}.

\bibitem{Levine1983}
M.~Levine,
\emph{Pluri-canonical divisors on Kähler manifolds},
Invent. Math. \textbf{74} (1983), no.~2, 293--303.
\href{https://doi.org/10.1007/BF01394318}{doi:10.1007/BF01394318}.

\bibitem{Levine1985}
M.~Levine,
\emph{Pluri-canonical divisors on Kähler manifolds II},
Duke Math. J. \textbf{52} (1985), no.~1, 61--65.
\href{https://doi.org/10.1215/S0012-7094-85-05205-6}{doi:10.1215/S0012-7094-85-05205-6}.

\bibitem{LiRaoWang2026}
M.-L.~Li, S.~Rao, and K.~Wang,
\emph{Deformation of nef adjoint canonical line bundles},
preprint, revised 27 October 2026.
\href{https://arxiv.org/abs/2510.23967v2}{arXiv:2510.23967v2}.

\bibitem{LPWZ2026}
Y.~Li, M.~Păun, Z.~Wang, and X.~Zhou,
\emph{Psh variation of singular Kähler--Einstein metrics and invariance
of plurigenera for smooth Kähler families},
preprint, 6 October 2026.
\href{https://arxiv.org/abs/2610.07826v1}{arXiv:2610.07826v1}.

\bibitem{LRW2019}
K.~Liu, S.~Rao, and X.~Wan,
\emph{Geometry of logarithmic forms and deformations of complex structures},
J. Algebraic Geom. \textbf{28} (2019), no.~4, 773--815.
\href{https://doi.org/10.1090/jag/723}{doi:10.1090/jag/723}.

\bibitem{LW2026}
K.~Liu and X.~Wan,
\emph{Invariance of logarithmic plurigenera and adjoint volumes in
Kähler families},
preprint, 19 October 2026.
\href{https://arxiv.org/abs/2610.26519v1}{arXiv:2610.26519v1}.

\bibitem{Mumford1966}
D.~Mumford,
\emph{Lectures on Curves on an Algebraic Surface},
Annals of Mathematics Studies, vol.~59,
Princeton University Press, Princeton, NJ, 1966.
\href{https://press.princeton.edu/books/paperback/9780691079936/lectures-on-curves-on-an-algebraic-surface}{Publisher's edition}.

\bibitem{Nakayama1986}
N.~Nakayama,
\emph{Invariance of the plurigenera of algebraic varieties under
minimal model conjectures},
Topology \textbf{25} (1986), no.~2, 237--251.
\href{https://doi.org/10.1016/0040-9383(86)90042-X}{doi:10.1016/0040-9383(86)90042-X}.

\bibitem{Nakayama1987}
N.~Nakayama,
\emph{The lower semi-continuity of the plurigenera of complex varieties},
in \emph{Algebraic Geometry, Sendai, 1985},
Adv. Stud. Pure Math., vol.~10,
North-Holland, Amsterdam, 1987, pp.~551--590.
\href{https://doi.org/10.2969/aspm/01010551}{doi:10.2969/aspm/01010551}.

\bibitem{Nakayama2004}
N.~Nakayama,
\emph{Zariski-decomposition and Abundance},
MSJ Memoirs, vol.~14, Mathematical Society of Japan, Tokyo, 2004.
\href{https://doi.org/10.2969/msjmemoirs/014010000}{doi:10.2969/msjmemoirs/014010000}.

\bibitem{OAI033}
OpenAI,
\emph{Orbifold and logarithmic Iitaka subadditivity},
preprint, 26 September 2026.
\href{https://github.com/openai/math/blob/fd4aeeb2ee4fc729c18d98444fed42fd0529eeeb/preprints/Orbifold-and-logarithmic-Iitaka-subadditivity-September-26-2026/paper.pdf}{Preprint}.

\bibitem{OAI034}
OpenAI,
\emph{Log abundance for compact Kähler spaces under logarithmic
Iitaka subadditivity},
preprint, 6 October 2026.
\href{https://github.com/openai/math/blob/fd4aeeb2ee4fc729c18d98444fed42fd0529eeeb/preprints/Log-abundance-for-compact-Kahler-spaces-under-logarithmic-Iitaka-subadditivity-October-6-2026/main.pdf}{Preprint}.

\bibitem{Ou2025}
W.~Ou,
\emph{A characterization of uniruled compact Kähler manifolds},
preprint, 30 January 2025.
\href{https://arxiv.org/abs/2501.18088v1}{arXiv:2501.18088v1}.

\bibitem{Paun2007}
M.~Păun,
\emph{Siu's invariance of plurigenera: a one-tower proof},
J. Differential Geom. \textbf{76} (2007), no.~3, 485--493.
\href{https://archive.ymsc.tsinghua.edu.cn/pacm_download/21/685-JDG-76-3-485-493.pdf}{Published article}.

\bibitem{RaoTsai2022}
S.~Rao and I.-H.~Tsai,
\emph{Invariance of plurigenera and Chow-type lemma},
Asian J. Math. \textbf{26} (2022), no.~4, 507--554.
\href{https://doi.org/10.4310/AJM.2022.v26.n4.a2}{doi:10.4310/AJM.2022.v26.n4.a2};
\href{https://arxiv.org/abs/2011.03306v2}{arXiv:2011.03306v2}.

\bibitem{Serre1956}
J.-P.~Serre,
\emph{Géométrie algébrique et géométrie analytique},
Ann. Inst. Fourier (Grenoble) \textbf{6} (1956), 1--42.
\href{https://doi.org/10.5802/aif.59}{doi:10.5802/aif.59}.

\bibitem{Siu1998}
Y.-T.~Siu,
\emph{Invariance of plurigenera},
Invent. Math. \textbf{134} (1998), no.~3, 661--673.
\href{https://doi.org/10.1007/s002220050276}{doi:10.1007/s002220050276}.

\bibitem{Siu2002}
Y.-T.~Siu,
\emph{Extension of twisted pluricanonical sections with plurisubharmonic
weight and invariance of semipositively twisted plurigenera for manifolds
not necessarily of general type},
in \emph{Complex Geometry} (Göttingen, 2000),
Springer, Berlin, 2002, pp.~223--277.
\href{https://doi.org/10.1007/978-3-642-56202-0_15}{doi:10.1007/978-3-642-56202-0\_15}.

\bibitem{Takayama2007}
S.~Takayama,
\emph{On the invariance and the lower semi-continuity of plurigenera
of algebraic varieties},
J. Algebraic Geom. \textbf{16} (2007), no.~1, 1--18.
\href{https://doi.org/10.1090/S1056-3911-06-00455-3}{doi:10.1090/S1056-3911-06-00455-3}.

\bibitem{TX2024}
N.~Tsakanikas and L.~Xie,
\emph{Remarks on the existence of minimal models of log canonical
generalized pairs},
Math. Z. \textbf{307} (2024), Article~20.
\href{https://doi.org/10.1007/s00209-024-03489-6}{doi:10.1007/s00209-024-03489-6};
\href{https://arxiv.org/abs/2301.09186v3}{arXiv:2301.09186v3}.

\end{thebibliography}

\end{document}
